% !TEX program = pdflatex
% !TEX encoding = UTF-8 Unicode

% ====================================================================
% ARCHIVO PRINCIPAL — Práctica Grupal Temas 2 y 3
% Asignatura: Dirección de Recursos Humanos I
% Doble Grado ADE-Ingenierías — Universidad de Granada
% Empresa analizada: CaixaBank
%
% INSTRUCCIONES DE COMPILACIÓN:
%   pdflatex practica.tex && pdflatex practica.tex
%   (dos pasadas para generar índice y referencias correctamente)
%
% Para la bibliografía, si añadís citas al archivo library.bib:
%   pdflatex practica.tex && bibtex practica && pdflatex practica.tex && pdflatex practica.tex
% ====================================================================

\documentclass[print, color]{ugrTFG}

% -------------------------------------------------------------------
%  VERSIÓN ELECTRÓNICA PARA TABLETA
%  Cambiad "print" por "tablet" para generar un PDF adaptado a tabletas.
% -------------------------------------------------------------------

% -------------------------------------------------------------------
%  INFORMACIÓN DEL TRABAJO Y LOS AUTORES
% -------------------------------------------------------------------

\renewcommand{\miTipoTrabajo}{Práctica Grupal}
\renewcommand{\miAsignatura}{Dirección de Recursos Humanos I}

\newcommand{\miTitulo}{Análisis y Diseño de Puestos de Trabajo y Planificación de Recursos Humanos en CaixaBank\xspace}

% Para un trabajo grupal, \miNombre contiene todos los autores.
% En la portada aparecerá como "Presentado por:" seguido de los nombres.
% Ajustadlos con vuestros nombres reales.
\newcommand{\miNombre}{%
  Nombre Apellidos 1\\
  Nombre Apellidos 2\\
  Nombre Apellidos 3\\
  Nombre Apellidos 4\\
  Nombre Apellidos 5%
}

\newcommand{\miGrado}{Doble Grado ADE-Ingenierías}
\newcommand{\miFacultad}{Facultad de Ciencias Empresariales y Escuela Técnica Superior de Ingenierías}
\newcommand{\miUniversidad}{Universidad de Granada}

% Nombre del profesor/a responsable de la asignatura
\newcommand{\miTutor}{
  Nombre del Profesor/a \\ \emph{Departamento de Organización de Empresas}
}

\newcommand{\miCurso}{2024-2025\xspace}

% -------------------------------------------------------------------
%  PAQUETES ADICIONALES PARA ESTE TRABAJO
% -------------------------------------------------------------------

\usepackage{tabularx}       % Tablas con columnas de anchura automática
\usepackage{multirow}       % Celdas que ocupan varias filas
\usepackage{float}          % Control de posición de figuras [H]
\usepackage{eurosym}        % Símbolo del euro: \euro{}
\usepackage{enumitem}       % Control avanzado de listas
\usepackage{tikz}           % Gráficos vectoriales (organigrama)
\usepackage{rotating}       % Figuras apaisadas (sidewaysfigure, sin problemas twoside)

% Cajas de color para recordatorios al equipo
\usepackage[most]{tcolorbox}
\tcbuselibrary{skins, breakable}

% Color corporativo UGR
\definecolor{ugrazul}{RGB}{0,75,135}

% Entorno "notaequipo": caja amarilla para recordatorios y preguntas guía.
% Usadlo cuando queráis dejar instrucciones a vuestros compañeros sobre
% qué información hay que buscar o qué preguntar en la entrevista.
\newtcolorbox{notaequipo}{
  enhanced,
  breakable,
  colback=yellow!15!white,
  colframe=orange!70!black,
  fonttitle=\sffamily\bfseries,
  title={Recordatorio del equipo},
  left=6pt, right=6pt, top=4pt, bottom=4pt,
  before skip=8pt, after skip=8pt
}

% -------------------------------------------------------------------
%  METADATOS PDF
% -------------------------------------------------------------------
\hypersetup{
  pdftitle={\miTitulo},
  pdfauthor={Grupo --- Dirección de Recursos Humanos I --- UGR},
  pdfsubject={Práctica Grupal — CaixaBank}
}

% ====================================================================
\begin{document}
% ====================================================================

\maketitle

% -------------------------------------------------------------------
%  FRONTMATTER
%  (páginas previas numeradas en romano: resumen, índice)
% -------------------------------------------------------------------
\frontmatter

% Resumen ejecutivo del trabajo
% !TeX root = ../practicat6t7t8.tex
% !TeX encoding = utf8

\chapter{Resumen}

El presente trabajo analiza las políticas de Formación y Desarrollo
(Tema~6), Evaluación y Gestión del Rendimiento (Tema~7) y Retribución
(Tema~8) aplicadas en CaixaBank, S.A., primer grupo bancario español
por volumen de activos, con más de 45.000 empleados.

En el ámbito de la formación, se examina el ciclo completo del programa
formativo de la entidad —identificación de necesidades, diseño,
implantación y evaluación—, así como el tiempo y coste asociados.
Se analizan igualmente las nuevas tendencias adoptadas: la plataforma
de \textit{e-learning} Virtaula, el modelo 70-20-10, la estructura de
universidades corporativas (escuelas especializadas internas),
el \textit{outdoor training} y la gamificación. Respecto a la gestión
de la carrera profesional, se describe el proceso de planificación
en sus cuatro fases —premisas, valoración, dirección y desarrollo—,
detallando las herramientas de rotación de puestos, \textit{coaching},
\textit{mentoring} e \textit{mentoring} inverso empleadas por la entidad,
junto con las tendencias de \textit{networking}, empleabilidad y marca
personal.

En el ámbito de la evaluación del rendimiento, se estudia la delimitación
de la evaluación (qué, cuándo y quién evalúa), los métodos objetivos,
subjetivos y mixtos utilizados, así como las estrategias de reforzamiento
positivo, negativo y la entrevista de evaluación. Se presta especial
atención a las tendencias de evaluación continua y por compromisos que
CaixaBank está implantando en sustitución de los sistemas anuales
tradicionales.

Finalmente, en materia retributiva, se analiza la estructura salarial de
cinco puestos representativos con referencia al Convenio Colectivo
sectorial de Banca, los tres tipos de equidad (interna, externa e
individual), el modelo híbrido de retribución por puesto y por
competencias, los principales incentivos financieros y elementos no
financieros, y las tendencias de retribución flexible y salario emocional.

La metodología combina el análisis de fuentes documentales públicas
—memoria anual, informe de sostenibilidad y Plan Estratégico
2025--2027 de CaixaBank— con una entrevista semiestructurada a una
directora de sucursal de la entidad.

\medskip

\textbf{Palabras clave:} formación corporativa, \textit{e-learning},
gestión de carrera profesional, \textit{coaching}, \textit{mentoring},
evaluación del rendimiento, retribución flexible, salario emocional,
CaixaBank, convenio colectivo de banca.

\cleardoublepage
\endinput


% Tabla de contenidos
% !TeX root = ../tfg.tex
% !TeX encoding = utf8
%
% ====================================================================
%  TABLA DE CONTENIDOS
% ====================================================================

\phantomsection
\pdfbookmark[0]{\contentsname}{toc}

\setcounter{tocdepth}{2}      % El índice muestra hasta subsecciones
\setcounter{secnumdepth}{3}   % Se numeran hasta subsubsecciones

\tableofcontents

% Listas de figuras y tablas (descomentad si las necesitáis)
% \phantomsection
% \listoffigures

% \phantomsection
% \listoftables

\cleardoublepage


% -------------------------------------------------------------------
%  MAINMATTER
%  (cuerpo del trabajo, numerado en arábigo)
% -------------------------------------------------------------------
\mainmatter

% ---- Capítulo 1: Introducción de la empresa ----
% !TeX root = ../tfg.tex
% !TeX encoding = utf8
%
% ====================================================================
%  CAPÍTULO 1 — INTRODUCCIÓN DE LA EMPRESA Y ORGANIGRAMA
%
%  RECORDATORIO AL EQUIPO:
%  Este capítulo sitúa al lector. Debe incluir:
%    - Descripción breve pero completa: sector, tamaño, historia reciente,
%      presencia geográfica, número de empleados, cifras clave.
%    - El organigrama: niveles jerárquicos, áreas funcionales y, si es
%      posible, el departamento de RRHH destacado.
%    - Cómo está estructurada la empresa formalmente vs. cómo funciona
%      realmente (funcional, divisional o matricial).
%    - Si tenéis el organigrama como imagen, guardadla en /img/
%      y usad \includegraphics. Si no, describidlo en texto.
%  Fuentes útiles: Memoria Anual CaixaBank, CNMV, web corporativa.
% ====================================================================

\chapter{La empresa: CaixaBank}

% ============================================================
\section{Descripción general}
% ============================================================

% TODO: Ampliad con datos reales obtenidos de la entrevista y fuentes
% oficiales (Memoria Anual CaixaBank, CNMV, web corporativa).
% Actualizad las cifras si han variado desde las que figuran aquí.

CaixaBank, S.A. es una entidad bancaria española con sede central en Valencia.
Surgida de la reorganización del Grupo ``la Caixa'', cotiza en el IBEX~35
y es actualmente uno de los mayores grupos financieros de España y Portugal
tras la fusión con Bankia completada en 2021.

\begin{table}[H]
  \centering
  \caption{Datos corporativos básicos de CaixaBank}
  \begin{tabularx}{\textwidth}{lX}
    \toprule
    \textbf{Parámetro}           & \textbf{Valor} \\
    \midrule
    Sector                       & Servicios financieros y bancarios \\
    Año de fundación             & 1904 (origen en Caja de Pensiones para la Vejez y de Ahorros) \\
    Sede central                 & Valencia, España \\
    Número de empleados          & $\sim$45.000 (grupo consolidado, 2024) \\
    Número de oficinas           & $\sim$4.000 en España \\
    Presencia internacional      & Portugal (BPI), Luxemburgo, Reino Unido, entre otros \\
    Cotización                   & IBEX~35 (ticker: CABK) \\
    Facturación / Margen bruto   & [Dato real a incluir] \\
    \bottomrule
  \end{tabularx}
\end{table}

% ============================================================
\section{Organigrama}
% ============================================================

% ------------------------------------------------------------------
%  COLORES PARA EL ORGANIGRAMA (tonos rojo plantilla UGR)
% ------------------------------------------------------------------
\definecolor{cCEO} {rgb}{0.40,0.02,0.02}
\definecolor{cL1}  {rgb}{0.58,0.07,0.07}
\definecolor{cGr}  {HTML}{155934}
\definecolor{cGr2} {HTML}{277048}
\definecolor{cEmp} {HTML}{7A8694}
\definecolor{cCard}{HTML}{FFFFFF}
\definecolor{cBord}{rgb}{0.80,0.68,0.68}
\definecolor{cConn}{rgb}{0.68,0.52,0.52}
\definecolor{cDark}{HTML}{1A0A0A}
\definecolor{cGray}{HTML}{5C4A4A}
\definecolor{cBg}  {rgb}{0.98,0.95,0.95}

% \nc{left_x}{top_y}{color}{titulo}{subtitulo}
\newcommand{\nc}[5]{%
  \fill[black!8,rounded corners=2.5pt](#1+.07,{-(#2)-.07})rectangle(#1+2.94+.07,{-(#2)-2.20-.07});
  \fill[cCard,rounded corners=2.5pt,draw=cBord,line width=.35pt](#1,{-(#2)})rectangle(#1+2.94,{-(#2)-2.20});
  \fill[#3,rounded corners=2.5pt](#1,{-(#2)})rectangle(#1+2.94,{-(#2)-.22});
  \fill[#3](#1,{-(#2)-.11})rectangle(#1+2.94,{-(#2)-.22});
  \fill[#3!22,draw=#3!55,line width=.4pt](#1+.26,{-(#2)-.58})circle(.22);
  \node[anchor=north west,text width=2.32cm,inner sep=0pt,align=left,
        font=\fontsize{5.5}{6.5}\bfseries\sffamily\selectfont,text=cDark]
    at(#1+.54,{-(#2)-.28}){#4};
  \def\tmpS{#5}\ifx\tmpS\empty\else
  \node[anchor=south west,text width=2.78cm,inner sep=0pt,align=left,
        font=\fontsize{4.6}{5.4}\sffamily\selectfont,text=cGray]
    at(#1+.08,{-(#2)-2.10}){#5};\fi}

\newcommand{\vc}[3]{\draw[cConn,line width=.5pt](#1,{-(#2)})--(#1,{-(#3)});}

\begin{sidewaysfigure}[p]
  \centering
  \scalebox{0.65}{%
  \begin{tikzpicture}[x=1cm,y=1cm,trim left=0cm]

    %% Fondo y barra título (A4 landscape = 29.7 x 21.0 cm)
    \fill[cBg](0,0)rectangle(29.7,-21.0);
    \fill[cCEO](0,0)rectangle(29.7,-.90);
    \node[text=white,font=\fontsize{10.5}{12}\bfseries\sffamily\selectfont,anchor=center]
      at(14.85,-.45){ORGANIGRAMA CORPORATIVO — CAIXABANK};

    %% CEO (centrado)
    \fill[black!8,rounded corners=3pt](12.518,-1.12)rectangle(15.918,-3.12);
    \fill[cCard,rounded corners=3pt,draw=cBord,line width=.45pt](12.448,-1.05)rectangle(15.848,-3.05);
    \fill[cCEO,rounded corners=3pt](12.448,-1.05)rectangle(15.848,-1.29);
    \fill[cCEO](12.448,-1.17)rectangle(15.848,-1.29);
    \fill[cCEO!22,draw=cCEO!55,line width=.5pt](12.748,-1.75)circle(.26);
    \node[anchor=north west,text width=2.86cm,inner sep=0pt,align=left,
          font=\fontsize{7}{8.5}\bfseries\sffamily\selectfont,text=cDark]
      at(13.118,-1.35){CONSEJERO DELEGADO};
    \node[anchor=south west,text width=3.22cm,inner sep=0pt,align=left,
          font=\fontsize{5.2}{6}\sffamily\selectfont,text=cGray]
      at(12.528,-2.95){18 / 38451 · 20 subordinados};

    %% Conectores CEO -> grilla
    \draw[cConn,line width=.6pt](14.148,-3.05)--(14.148,-3.30);
    \draw[cConn,line width=.6pt](1.47,-3.30)--(24.918,-3.30);
    \foreach \cx in {1.47,4.614,7.758,10.902,14.046,17.190,20.334,23.478}{
      \draw[cConn,line width=.55pt](\cx,-3.30)--(\cx,-3.45);}

    %% FILA 1
    \nc{0.000}{3.45}{cL1}{DIRECTOR GABINETE CONSEJ. DELEGADO}{5/5 · 5 subordinados de matriz}
    \nc{3.144}{3.45}{cL1}{DIRECTOR PERSONAS}{15/211 · 17 subordinados de...}
    \nc{6.288}{3.45}{cL1}{MANAGING DIRECTOR DESARROLLO CORPORATIVO}{8/15 · 8 subordinados de mat...}
    \nc{9.432}{3.45}{cL1}{DIRECTOR SEGUROS}{2/16 · 2 subordinados de mat...}
    \nc{12.576}{3.45}{cEmp}{EMPLEADA}{}
    \nc{15.720}{3.45}{cL1}{DIRECTOR PERSONAS}{15/211 · 17 subordinados de...}
    \nc{18.864}{3.45}{cL1}{DIRECTOR SEGUROS}{2/16 · 2 subordinados de mat...}
    \nc{22.008}{3.45}{cGr}{DIRECTOR NEGOCIO}{30/31327 · 32 subordinados...}
    %% FILA 2
    \nc{0.000}{5.87}{cL1}{DIRECTOR CORPORATE \& INVESTMENT BANKING}{22/917 · 22 subordinados de...}
    \nc{3.144}{5.87}{cL1}{DIRECTOR FINANCIERA}{9/190 · 9 subordinados de m...}
    \nc{6.288}{5.87}{cEmp}{EMPLEADA}{}
    \nc{9.432}{5.87}{cL1}{DIRECTORA COMUNICACION Y RELAC. INSTITUCIONALES}{12/106 · 10 subordinados de...}
    \nc{12.576}{5.87}{cL1}{DIRECTORA TRANSFORM. DIGITAL Y ADVANCED ANALYTICS}{13/838 · 16 subordinados de...}
    \nc{15.720}{5.87}{cEmp}{EMPLEADA}{}
    \nc{18.864}{5.87}{cGr}{DIRECTOR NEGOCIO}{30/31327 · 32 subordinados...}
    \nc{22.008}{5.87}{cGr}{DIRECTOR TERRITORIAL ANDALUCIA}{28/4250 · 28 subordinados d...}
    %% FILA 3
    \nc{0.000}{8.29}{cL1}{DIRECTOR MEDIOS}{18/2237 · 19 subordinados d...}
    \nc{3.144}{8.29}{cL1}{DIRECTOR ADM. Y SEGUIMIENTO DEL RIESGO DE CREDITO}{9/1430 · 9 subordinados de...}
    \nc{6.288}{8.29}{cL1}{DIRECTOR CORPORATE \& INVESTMENT BANKING}{22/917 · 22 subordinados de...}
    \nc{9.432}{8.29}{cL1}{DIRECTORA GESTION DE RIESGOS Y CUMPLIMIENTO}{11/482 · 11 subordinados de...}
    \nc{12.576}{8.29}{cL1}{SECRETARIO GENERAL}{9/268 · 9 subordinados de m...}
    \nc{15.720}{8.29}{cL1}{DIRECTORA TRANSFORM. DIGITAL Y ADVANCED ANALYTICS}{13/838 · 16 subordinados de...}
    \nc{18.864}{8.29}{cL1}{DIRECTOR ADM. Y SEGUIMIENTO DEL RIESGO DE CREDITO}{9/1430 · 9 subordinados de...}
    \nc{22.008}{8.29}{cGr}{DIRECTOR COMERCIAL RED GRANADA -- JAEN}{9/925 · 9 subordinados de m...}
    \nc{25.152}{8.29}{cGr2}{SUBDIRECTOR COMERCIAL}{}
    %% FILA 4
    \nc{0.000}{10.71}{cL1}{DIRECTOR CONTABILIDAD, CTRL DE GESTION Y CAPITAL}{15/391 · 13 subordinados de...}
    \nc{3.144}{10.71}{cEmp}{EMPLEADA}{}
    \nc{6.288}{10.71}{cL1}{DIRECTOR PUBLIC AFFAIRS}{10/10 · 10 subordinados de...}
    \nc{9.432}{10.71}{cL1}{DIRECTOR PERSONAS}{15/211 · 17 subordinados de...}
    \nc{12.576}{10.71}{cL1}{DIRECTOR FINANCIERA}{9/190 · 9 subordinados de m...}
    \nc{15.720}{10.71}{cL1}{DIRECTORA FINANCIERA}{9/190 · 9 subordinados de m...}
    \nc{18.864}{10.71}{cEmp}{EMPLEADA}{}
    \nc{22.008}{10.71}{cGr}{DIRECTOR AREA DE NEGOCIO}{24/95 · 24 subordinados de...}
    \nc{25.152}{10.71}{cGr2}{SUBDIRECTOR COMERCIAL}{}
    %% FILA 5
    \nc{0.000}{13.13}{cEmp}{EMPLEADA}{}
    \nc{3.144}{13.13}{cL1}{DIRECTOR GABINETE CONSEJ. DELEGADO}{5/5 · 5 subordinados de matriz}
    \nc{6.288}{13.13}{cEmp}{EMPLEADA}{}
    \nc{9.432}{13.13}{cL1}{DIRECTORA COMUNICACION Y RELAC. INSTITUCIONALES}{12/106 · 10 subordinados de...}
    \nc{12.576}{13.13}{cEmp}{EMPLEADA}{}
    \nc{15.720}{13.13}{cL1}{DIRECTORA TRANSFORM. DIGITAL Y ADVANCED ANALYTICS}{13/838 · 16 subordinados de...}
    \nc{18.864}{13.13}{cL1}{DIRECTORA FINANCIERA}{9/190 · 9 subordinados de m...}
    \nc{22.008}{13.13}{cGr}{DIRECTORA}{2/2 · 2 subordinados de matriz}
    %% FILA 6
    \nc{3.144}{15.55}{cL1}{DIRECTOR SEGUROS}{2/16 · 2 subordinados de mat...}
    \nc{6.288}{15.55}{cL1}{DIRECTOR ADM. Y SEGUIMIENTO DEL RIESGO DE CREDITO}{9/1430 · 9 subordinados de...}
    \nc{9.432}{15.55}{cL1}{DIRECTORA GESTION DE RIESGOS Y CUMPLIMIENTO}{11/482 · 11 subordinados de...}
    \nc{12.576}{15.55}{cL1}{SECRETARIO GENERAL}{9/268 · 9 subordinados de m...}
    %% FILA 7
    \nc{3.144}{17.97}{cGr}{DIRECTOR NEGOCIO}{30/31327 · 32 subordinados...}
    \nc{6.288}{17.97}{cEmp}{EMPLEADA}{}

    %% CONECTORES VERTICALES
    \vc{1.470}{5.65}{5.87}\vc{1.470}{8.07}{8.29}\vc{1.470}{10.49}{10.71}\vc{1.470}{12.91}{13.13}
    \vc{4.614}{5.65}{5.87}\vc{4.614}{8.07}{8.29}\vc{4.614}{10.49}{10.71}\vc{4.614}{12.91}{13.13}\vc{4.614}{15.33}{15.55}\vc{4.614}{17.75}{17.97}
    \vc{7.758}{5.65}{5.87}\vc{7.758}{8.07}{8.29}\vc{7.758}{10.49}{10.71}\vc{7.758}{12.91}{13.13}\vc{7.758}{15.33}{15.55}\vc{7.758}{17.75}{17.97}
    \vc{10.902}{5.65}{5.87}\vc{10.902}{8.07}{8.29}\vc{10.902}{10.49}{10.71}\vc{10.902}{12.91}{13.13}\vc{10.902}{15.33}{15.55}
    \vc{14.046}{5.65}{5.87}\vc{14.046}{8.07}{8.29}\vc{14.046}{10.49}{10.71}\vc{14.046}{12.91}{13.13}\vc{14.046}{15.33}{15.55}
    \vc{17.190}{5.65}{5.87}\vc{17.190}{8.07}{8.29}\vc{17.190}{10.49}{10.71}\vc{17.190}{12.91}{13.13}
    \vc{20.334}{5.65}{5.87}\vc{20.334}{8.07}{8.29}\vc{20.334}{10.49}{10.71}\vc{20.334}{12.91}{13.13}
    \vc{23.478}{5.65}{5.87}\vc{23.478}{8.07}{8.29}\vc{23.478}{10.49}{10.71}\vc{23.478}{12.91}{13.13}
    %% Conector col7->col8 (subdirectores comerciales)
    \draw[cConn,line width=.5pt](24.948,-9.490)--(25.152,-9.490);
    \draw[cConn,line width=.5pt](24.948,-11.910)--(25.152,-11.910);
    \draw[cConn,line width=.5pt](24.948,-9.490)--(24.948,-11.910);

  \end{tikzpicture}%
  }% fin \resizebox
  \caption{Organigrama corporativo de CaixaBank: Consejero Delegado y principales directivos de primer y segundo nivel (fuente: Interna\,/\,elaboración propia, 2026).
           \textcolor{cL1}{$\blacksquare$}~Direcciones corporativas de sede;\quad
           \textcolor{cGr}{$\blacksquare$}~Red territorial de negocio;\quad
           \textcolor{cEmp}{$\blacksquare$}~Empleado/a sin equipo a cargo.}
  \label{fig:organigrama}
\end{sidewaysfigure}

% ============================================================
\section{El departamento de Personas y Organización (RRHH)}
% ============================================================
A continuación, se muestra el organigrma del departamento:


% ---- Capítulo 2: Tema 2 — Análisis y Diseño de Puestos ----
% !TeX root = ../tfg.tex
% !TeX encoding = utf8
%
% ====================================================================
%  CAPÍTULO 2 — TEMA 2: ANÁLISIS Y DISEÑO DE PUESTOS DE TRABAJO
%
%  Este capítulo cubre todos los apartados del Tema 2 de la asignatura:
%    2.1  Análisis de los puestos de trabajo (APT): etapas en CaixaBank
%    2.2  Descripción y especificación de tres puestos
%    2.3  Formas de rediseño de puestos
%    2.4  Nuevas tendencias: ADP por competencias
% ====================================================================

\chapter{Tema 2: Análisis y Diseño de Puestos de Trabajo}


% ================================================================
\section{Análisis de los puestos de trabajo en CaixaBank}
% ================================================================

% RECORDATORIO AL EQUIPO:
% Las etapas del análisis de puestos (APT) son generalmente:
%   1. Planificación del APT (objetivos, recursos)
%   2. Preparación (diseño de instrumentos: cuestionarios, entrevistas...)
%   3. Recogida de información (observación, entrevista, cuestionario, diario...)
%   4. Análisis y tratamiento de datos
%   5. Elaboración de la descripción y especificación del puesto
%   6. Actualización y mantenimiento
%
% PARA CADA ETAPA INDICAD:
%   - ¿La realiza CaixaBank? ¿Cómo exactamente?
%   - Ventajas e inconvenientes de su enfoque
%   - ¿Está la empresa satisfecha con ese proceso?
%   - ¿Habría maneras más eficientes de hacerlo?

El Análisis del Puesto de Trabajo (APT) es el proceso a través del cual la empresa recopila y analiza la información sobre el puesto de trabajo con la intención de identificar las tareas, obligaciones y sus responsabilidades, de forma que sirvan para establecer el perfil de persona que debería ocuparlo.

A continuación se detalla en qué medida CaixaBank desarrolla cada una
de las etapas clásicas del APT:

% ---- TABLA RESUMEN DE ETAPAS ----
\begin{table}[H]
  \centering
  \small
  \caption{Resumen de etapas del APT en CaixaBank}
  \begin{tabularx}{\textwidth}{p{0.30\textwidth} p{0.12\textwidth} X}
    \toprule
    \textbf{Etapa} & \textbf{¿Se realiza?} & \textbf{Observaciones} \\
    \midrule
    1. Planificación del APT          & Sí & (Liderado por Secretaría Técnica de Negocios y Deloitte). \\
    2. Preparación de instrumentos    & Sí & (Enfoque bottom-up, basado en las necesidades del cliente). \\
    3. Recogida de información        & Sí & (Adaptada a la ubicación y contexto de la oficina). \\
    4. Análisis y tratamiento         & Sí & (Apoyado por consultoría externa). \\
    5. Descripción y especificación   & Sí & (Totalmente orientado a competencias y soft skills). \\
    6. Actualización y mantenimiento  & Sí & (Dinámico y ajustado a la demanda del mercado). \\
    \bottomrule
  \end{tabularx}
\end{table}

% ---- ETAPA 1 ----
\subsection{Etapa 1 — Planificación del análisis de puestos}

% TODO: ¿Planifica formalmente CaixaBank el APT? ¿Tiene un departamento
% o equipo dedicado? ¿Con qué periodicidad lo actualiza?
% ¿Lo hace internamente o externaliza (consultoría)?

A diferencia de los modelos teóricos tradicionales, en CaixaBank la planificación del análisis de puestos no recae de forma exclusiva en el departamento de Recursos Humanos. Esta etapa está externalizada y liderada por la Secretaría Técnica de Negocios, utilizando un modelo desarrollado por la consultora externa Deloitte.

\textbf{Ventajas del enfoque adoptado:}
\begin{itemize}
  \item Alineación con el negocio: Al ser diseñado por la Secretaría Técnica, el puesto nace con un enfoque 100\% comercial y estratégico, garantizando que el empleado aporte valor real a la cuenta de resultados de la entidad.
  \item Objetividad y profesionalidad: Contar con una consultora experta (Deloitte) aporta metodologías probadas en el mercado y evita sesgos internos.
\end{itemize}

\textbf{Inconvenientes / áreas de mejora:}
\begin{itemize}
  \item Coste económico: Externalizar este proceso a consultoras de primer nivel supone un alto coste para la organización.
  \item Distanciamiento de RRHH: Que RRHH no defina el puesto puede generar desajustes a la hora de realizar la posterior selección o formación si la comunicación con la Secretaría Técnica no es fluida.
\end{itemize}

% ---- ETAPA 2 ----
\subsection{Etapa 2 — Preparación de los instrumentos de recogida}

% TODO: ¿Qué métodos de recogida de información utiliza CaixaBank?
% Métodos habituales: cuestionario estructurado, entrevista al ocupante,
% entrevista al supervisor, observación directa, diario de actividades,
% incidentes críticos, conferencia técnica de expertos.
% Indicad cuáles SÍ usa, cuáles NO y por qué (coste, tipo de puesto...).

CaixaBank utiliza un enfoque radicalmente distinto al tradicional: lo hacen "de abajo arriba" (bottom-up). En lugar de que la dirección imponga unas tareas cerradas, la recogida de información se realiza partiendo de las necesidades reales del cliente en la base.

Dado que el diseño lo elabora Negocio junto a Deloitte, se deduce que el método principal para recopilar datos no es la simple observación, sino los grupos de expertos. Este método es muy útil cuando lo que se pretende es el diseño de un nuevo puesto o un rediseño sustancial, reuniendo a directores de zona o especialistas para definir qué se necesita exactamente en la red comercial.

% ---- ETAPA 3 ----
\subsection{Etapa 3 — Recogida de información}

% TODO: ¿Quién recoge la información? ¿El propio empleado, su supervisor,
% un analista de RRHH? ¿Hay sesgo? ¿Cómo garantizan la fiabilidad?

Además, la recogida de información tiene muy en cuenta el contexto local, ya que el diseño del puesto y las habilidades requeridas varían "dependiendo del lugar donde está la oficina".

% ---- ETAPAS 4, 5 y 6 ----
\subsection{Etapas 4, 5 y 6 — Análisis, descripción y actualización}

% TODO: ¿Existe un catálogo formal de puestos? ¿Con qué frecuencia se revisa?
% En un banco como CaixaBank la aparición de nuevas figuras (analista de datos,
% ciberseguridad, UX bancario) exige actualizaciones frecuentes.
% ¿Lo hace CaixaBank? ¿Cómo gestionan la obsolescencia de los perfiles?

En la fase de descripción y especificación del puesto, CaixaBank aplica directamente las nuevas tendencias en RRHH: el análisis de puestos basado en competencias.

El modelo de Deloitte no redacta descripciones de puestos basadas únicamente en un listado rígido de tareas repetitivas, sino que implica la redacción de la descripción teniendo en cuenta los conocimientos, experiencias y, sobre todo, las competencias que se requieren para que se pueda desempeñar el puesto con éxito.

Le dan una importancia capital a las soft skills (competencias blandas como la empatía, la capacidad de negociación, la adaptabilidad o el trato al cliente). Esto permite que los puestos no se queden obsoletos rápidamente, ya que si cambian las herramientas informáticas del banco, el empleado con las soft skills adecuadas (flexibilidad, aprendizaje) se adaptará sin necesidad de tener que rehacer por completo la descripción formal de su puesto.


% ================================================================
\section{Descripción y especificación de tres puestos de CaixaBank}
% ================================================================

% RECORDATORIO AL EQUIPO:
% Hay que presentar TRES puestos distintos de la empresa.
% Para cada puesto debéis elaborar:
%
% A) DESCRIPCIÓN DEL PUESTO (el QUÉ se hace):
%    - Identificación: denominación, departamento, nivel jerárquico,
%      superior inmediato, subordinados, fecha de elaboración.
%    - Resumen del puesto (misión en 2-3 líneas).
%    - Funciones y tareas principales (ordenadas por importancia/frecuencia).
%    - Responsabilidades: personas, equipos, presupuesto, información.
%    - Condiciones de trabajo: horario, lugar, movilidad, riesgos.
%    - Relaciones internas y externas.
%
% B) ESPECIFICACIÓN DEL PUESTO (el QUIÉN lo puede hacer):
%    - Formación mínima requerida.
%    - Experiencia mínima.
%    - Conocimientos específicos.
%    - Competencias / habilidades (técnicas y conductuales).
%    - Rasgos de personalidad relevantes.
%
% PARA CADA PUESTO VALORAD:
%    - ¿La descripción está actualizada? ¿Refleja la realidad?
%    - ¿La especificación es adecuada (ni sobre- ni infracualificación)?
%    - ¿Está la empresa satisfecha con estas fichas?
%    - ¿Podría mejorarse? ¿Cómo?
%
% Elegid tres puestos de diferentes niveles jerárquicos o áreas:
%   - Director/a de Oficina (nivel directivo / mando intermedio)
%   - Gestor/a de Banca Privada (nivel técnico especializado)
%   - Asesor/a Comercial de Oficina (nivel operativo)


% ============================================================
\subsection{Puesto 1: Director/a de Oficina}
% ============================================================

\subsubsection{Descripción del puesto}

\begin{table}[H]
  \centering
  \small
  \caption{Ficha de identificación — Puesto 1}
  \begin{tabularx}{\textwidth}{lX}
    \toprule
    \textbf{Campo}                & \textbf{Contenido} \\
    \midrule
    Denominación del puesto       & Director/a de Oficina \\
    Código / referencia interna   & [Si existe] \\
    Departamento / Área           & Red Territorial / Oficinas \\
    Nivel jerárquico              & Mando intermedio \\
    Superior inmediato            & Director/a Territorial de Zona \\
    Número de subordinados        & Entre 3 y 10 personas según el tamaño de la Store u oficina \\
  \end{tabularx}
\end{table}

\textbf{Misión del puesto:}

Gestionar, coordinar y liderar el equipo de la oficina bancaria, asegurando el cumplimiento de los objetivos comerciales, la rentabilidad de la sucursal y la excelencia en la satisfacción del cliente, operando siempre bajo los estrictos estándares de riesgo y cumplimiento normativo del Grupo CaixaBank.

\textbf{Funciones y tareas principales:}
\begin{enumerate}
  \item Planificar y asignar los objetivos comerciales individuales al equipo de la oficina.
  \item Realizar el seguimiento sistemático de la actividad comercial y proponer planes de acción ante desviaciones.
  \item Gestionar directamente la cartera de clientes de mayor valor o riesgo de la sucursal.
  \item Liderar el desarrollo profesional, la motivación y la evaluación del desempeño de los empleados a su cargo.
  \item Garantizar el cumplimiento de las normativas de auditoría interna, prevención de blanqueo de capitales y MiFID II.
\end{enumerate}

\textbf{Responsabilidades:}
\begin{itemize}
  \item \textit{Sobre personas:} Liderazgo, supervisión, evaluación y motivación de un equipo (habitualmente entre 3 y 10 personas según el tamaño de la Store u oficina).
  \item \textit{Sobre equipos / materiales:} Gestión de la infraestructura de la oficina y seguridad física.
  \item \textit{Sobre información confidencial:} Acceso a datos patrimoniales de clientes, gestión de riesgos crediticios y morosidad.
  \item \textit{Sobre presupuesto:} Responsabilidad directa sobre la cuenta de resultados (P\&L) de la oficina y el cumplimiento del presupuesto anual asignado por la Dirección Territorial.
\end{itemize}

\textbf{Condiciones de trabajo:}
\begin{itemize}
  \item Horario: Jornada completa, habitualmente horario partido (mañana y tarde, especialmente en el modelo "Oficina Store" de CaixaBank).
  \item Lugar: Oficina física, con escasa posibilidad de teletrabajo.
  \item Movilidad / desplazamientos: Sí, ocasionales para visitas a empresas clientes o reuniones en la Dirección Territorial.
  \item Riesgos laborales: Alto nivel de estrés por presión comercial y cumplimiento de objetivos; riesgos psicosociales y ergonomía por uso de pantallas.
\end{itemize}

\subsubsection{Especificación del puesto}

\begin{table}[H]
  \centering
  \small
  \caption{Requisitos del candidato — Puesto 1}
  \begin{tabularx}{\textwidth}{lX}
    \toprule
    \textbf{Requisito}        & \textbf{Descripción} \\
    \midrule
    Formación mínima          & Grado en ADE, Economía, Derecho o similar \\
    Formación complementaria  & Certificación MiFID II (imprescindible) y Ley de Contrato de Crédito Inmobiliario (LCCI) \\
    Experiencia mínima        & Al menos 5 años en el sector bancario, con un mínimo de 2-3 años de experiencia previa en gestión de equipos o como Subdirector/a \\
    Conocimientos técnicos    & Análisis de riesgos de empresas y particulares, cuenta de resultados, productos de inversión y financiación \\
    Idiomas                   & Español nativo; nivel alto del idioma cooficial (si aplica a la CCAA); inglés nivel B2 valorable \\
    Competencias conductuales & Liderazgo transformacional, toma de decisiones bajo presión, alta orientación a resultados y resiliencia \\
    Herramientas / software   & Terminal Financiero propio de CaixaBank, CRM interno, Office 365 \\
    \bottomrule
  \end{tabularx}
\end{table}

En general, la especificación interna del puesto de Director/a de Oficina es más exigente que las ofertas externas en:
\begin{itemize}
  \item \textbf{Reclutamiento:} las ofertas externas tienden a centrarse en habilidades comerciales y cumplimiento de KPI; la ficha interna refuerza el enfoque en riesgo y cumplimiento.
  \item \textbf{Requisitos formativos:} el grado y la certificación MiFID II son estándar, pero la inclusión de LCCI y experiencia en gestión de equipos puede reducir el pool de candidatos.
  \item \textbf{Competencias:} la ficha interna prioriza resiliencia y liderazgo; las ofertas externas suelen insistir más en orientación comercial y resultados.
\end{itemize}

% ============================================================
\subsection{Puesto 2: Gestor/a de Banca Privada}
% ============================================================

\subsubsection{Descripción del puesto}

\begin{table}[H]
  \centering
  \small
  \caption{Ficha de identificación — Puesto 2}
  \begin{tabularx}{\textwidth}{lX}
    \toprule
    \textbf{Campo}                & \textbf{Contenido} \\
    \midrule
    Denominación del puesto       & Gestor/a de Banca Privada \\
    Departamento / Área           & Banca Privada / CaixaBank Private Banking \\
    Nivel jerárquico              & Técnico especializado \\
    Superior inmediato            & Director/a de Banca Privada \\
    Número de subordinados        & Ninguno / gestión de cartera propia \\
    \bottomrule
  \end{tabularx}
\end{table}

\textbf{Misión del puesto:}

Proporcionar un asesoramiento patrimonial y financiero integral y altamente personalizado a clientes del segmento de rentas altas (Banca Privada), buscando la optimización fiscal y la rentabilidad de sus inversiones para fidelizar e incrementar los activos bajo gestión (AUM).

\textbf{Funciones y tareas principales:}
\begin{enumerate}
  \item Diseñar, presentar y ejecutar propuestas de inversión personalizadas acordes al perfil de riesgo del cliente.
  \item Realizar un seguimiento proactivo de los mercados financieros y rebalancear las carteras de los clientes cuando sea necesario.
  \item Captar nuevo patrimonio y nuevos clientes dentro del segmento objetivo.
  \item Coordinar con especialistas internos (fiscalistas, gestores de fondos) la estructuración del patrimonio del cliente.
\end{enumerate}

\textbf{Responsabilidades:}
\begin{itemize}
  \item \textit{Sobre personas:} Ninguna. Es un perfil de contribuidor individual.
  \item \textit{Sobre equipos / materiales:} Equipos informáticos corporativos en movilidad (portátil, smartphone).
  \item \textit{Sobre información confidencial:} Confidencialidad extrema sobre grandes patrimonios, estructuras societarias y datos sensibles.
  \item \textit{Sobre presupuesto:} Responsabilidad sobre el crecimiento del volumen de negocio (Assets Under Management) de su cartera asignada.
\end{itemize}

\textbf{Condiciones de trabajo:}
\begin{itemize}
  \item Horario: Jornada completa con alta flexibilidad horaria para adaptarse a las necesidades de clientes VIP.
  \item Lugar: Centro de Banca Privada y alta movilidad.
  \item Movilidad / desplazamientos: Muy frecuente (visitas a domicilios, empresas de clientes o eventos de networking).
  \item Riesgos laborales: Estrés ligado a la volatilidad de los mercados financieros.
\end{itemize}

\subsubsection{Especificación del puesto}

\begin{table}[H]
  \centering
  \small
  \caption{Requisitos del candidato — Puesto 2}
  \begin{tabularx}{\textwidth}{lX}
    \toprule
    \textbf{Requisito}        & \textbf{Descripción} \\
    \midrule
    Formación mínima          & Grado en ADE, Economía o Ciencias Actuariales \\
    Certificaciones           & Certificación de Asesor Financiero (Nivel EFPA o equivalente) obligatoria; CFA muy valorable \\
    Experiencia mínima        & Entre 3 y 5 años de experiencia demostrable en gestión de patrimonios o mercados financieros \\
    Conocimientos técnicos    & Profundo conocimiento macroeconómico, fiscalidad de inversiones, SICAVs, fondos de inversión alternativos \\
    Competencias conductuales & Excelentes habilidades de comunicación, empatía, discreción absoluta, proactividad y tolerancia a la frustración \\
    \bottomrule
  \end{tabularx}
\end{table}


En comparación con el puesto de Director/a de Oficina, el Gestor/a de Banca Privada muestra:
\begin{itemize}
  \item \textbf{Mayor especialización técnica:} trabaja con productos complejos (SICAVs, fondos alternativos, fiscalidad patrimonial) y requiere certificaciones como EFPA/CFA.
  \item \textbf{Relación cliente más exclusiva:} gestiona carteras de clientes VIP y ofrece soluciones a medida, mientras que el director combina atención cliente con gestión de oficina.
  \item \textbf{Responsabilidad individual:} no supervisa equipo ni P\&L de la oficina, pero tiene responsabilidad directa sobre el crecimiento de los activos bajo gestión.
\end{itemize}

\textbf{Conclusión:} CaixaBank distingue claramente entre banca personal (oficina) y banca privada, manteniendo roles con distintos niveles técnicos y tipos de cliente, aunque ambos mantienen el foco en la venta y en la gestión del riesgo.


% ============================================================
\subsection{Puesto 3: Asesor/a Comercial de Oficina}
% ============================================================

\subsubsection{Descripción del puesto}

\begin{table}[H]
  \centering
  \small
  \caption{Ficha de identificación — Puesto 3}
  \begin{tabularx}{\textwidth}{lX}
    \toprule
    \textbf{Campo}                & \textbf{Contenido} \\
    \midrule
    Denominación del puesto       & Asesor/a Comercial de Oficina \\
    Departamento / Área           & Red Comercial \\
    Nivel jerárquico              & Operativo / técnico base \\
    Superior inmediato            & Director/a de Oficina \\
    Número de subordinados        & Ninguno \\
    Fecha de elaboración / revisión & [Fecha] \\
    \bottomrule
  \end{tabularx}
\end{table}

\textbf{Misión del puesto:}

Atender de forma proactiva y excelente a los clientes en la sucursal, resolviendo sus operativas diarias y detectando oportunidades para comercializar productos y servicios financieros y aseguradores básicos, derivando al cliente a perfiles especializados cuando sea necesario.

\textbf{Funciones y tareas principales:}
\begin{enumerate}
  \item Recepción, atención y derivación de clientes en la sucursal (labor de "filtro" o bienvenida).
  \item Venta cruzada de productos transaccionales, tarjetas, seguros de auto/hogar y préstamos al consumo.
  \item Resolución de incidencias operativas y soporte en el uso de los canales digitales (cajeros, app CaixaBankNow).
  \item Realización de campañas comerciales telefónicas (telemarketing) a la cartera de clientes asignada.
\end{enumerate}

\textbf{Responsabilidades:}
\begin{itemize}
  \item \textit{Sobre personas:} Ninguna.
  \item \textit{Sobre equipos / materiales:} Terminal financiero.
  \item \textit{Sobre información confidencial:} Protección de datos personales (RGPD) de los clientes de la oficina.
  \item \textit{Sobre presupuesto:} Cumplimiento de objetivos comerciales individuales (número de tarjetas, pólizas de seguro, etc.).
\end{itemize}

\textbf{Condiciones de trabajo:}
\begin{itemize}
  \item Horario: Jornada continua (normalmente de 8:00 a 15:00) o partida en oficinas Store, dependiendo del contrato.
  \item Lugar: Puesto fijo en la zona de atención al público de la sucursal.
  \item Movilidad / desplazamientos: Nula o muy baja.
  \item Riesgos laborales: Fatiga vocal, exposición al público (riesgo de clientes conflictivos), sedentarismo.
\end{itemize}

\subsubsection{Especificación del puesto}

\begin{table}[H]
  \centering
  \small
  \caption{Requisitos del candidato — Puesto 3}
  \begin{tabularx}{\textwidth}{lX}
    \toprule
    \textbf{Requisito}        & \textbf{Descripción} \\
    \midrule
    Formación mínima          & Grado Universitario (preferiblemente en rama económica/empresarial) o FP de Grado Superior en Administración y Finanzas \\
    Certificaciones           & Certificación MiFID II (muchas veces el banco asume esta formación tras la contratación) \\
    Experiencia mínima        & 0 a 2 años (es un puesto habitual de entrada o junior) \\
    Conocimientos técnicos    & Conocimiento de productos financieros básicos, ofimática y agilidad en sistemas informáticos \\
    Competencias conductuales & Gran orientación al cliente, dotes comerciales, empatía, paciencia y capacidad de trabajo en equipo \\
    \bottomrule
  \end{tabularx}
\end{table}


  \textbf{Rotación y ajuste del perfil:}
  \begin{itemize}
    \item \textbf{Rotación alta en nivel junior:} suele existir una brecha entre las expectativas de desarrollo y la realidad operativa; se puede mejorar con formación y acompañamiento inicial más estructurado.
    \item \textbf{Dificultad para cubrir vacantes:} las ofertas pueden describir un puesto más amplio del que luego resulta ser en el día a día; esto alarga el proceso de selección y genera fricción.
    \item \textbf{Desajuste perfil-puesto:} si se prioriza actitud comercial en detrimento de competencias técnicas básicas, los nuevos empleados pueden frustrarse rápidamente.
  \end{itemize}

  \textbf{Sugerencia:} Fortalecer el onboarding con formación práctica, clarificar objetivos comerciales desde el inicio y revisar incentivos para mejorar la retención.



% ================================================================
\section{Formas de rediseño de puestos en CaixaBank}
% ================================================================

% RECORDATORIO AL EQUIPO:
% Las principales formas de rediseño de puestos son:
%
%   1. ROTACIÓN DE PUESTOS (job rotation): los empleados rotan entre puestos.
%      Ventajas: polivalencia, motivación, reducción de la monotonía.
%      Inconvenientes: costes de formación, pérdida de especialización.
%
%   2. AMPLIACIÓN DE TAREAS (job enlargement / ensanchamiento horizontal):
%      se añaden tareas del mismo nivel de complejidad.
%      Ventajas: reduce la monotonía, mayor variedad.
%      Inconvenientes: puede aumentar la carga sin añadir motivación real.
%
%   3. ENRIQUECIMIENTO DE PUESTOS (job enrichment / Hackman & Oldham):
%      se añaden responsabilidades de mayor complejidad (carga vertical).
%      Las cinco características centrales del modelo:
%        - Variedad de habilidades
%        - Identidad de la tarea
%        - Significado de la tarea
%        - Autonomía
%        - Retroalimentación
%      Ventajas: mayor motivación intrínseca, satisfacción, compromiso.
%      Inconvenientes: requiere empleados que quieran / puedan asumir más.
%
%   4. EQUIPOS DE TRABAJO AUTOGESTIONADOS (self-managed teams).
%   5. TELETRABAJO / TRABAJO FLEXIBLE (flexibilización del puesto).
%   6. JOB CRAFTING: el empleado adapta el puesto a sus fortalezas
%      de forma proactiva (tendencia emergente).
%
% PARA CADA FORMA INDICAD:
%   - ¿La aplica CaixaBank? ¿En qué puestos / áreas?
%   - Ventajas e inconvenientes concretos observados.
%   - ¿Está la empresa satisfecha?
%   - ¿Qué formas serían interesantes aplicar y por qué?

\subsection{Formas de rediseño aplicadas actualmente}

\begin{table}[H]
  \centering
  \small
  \caption{Aplicación de formas de rediseño de puestos en CaixaBank}
  \begin{tabularx}{\textwidth}{p{0.30\textwidth} p{0.12\textwidth} X}
    \toprule
    \textbf{Forma de rediseño}       & \textbf{¿Se aplica?} & \textbf{Descripción / observaciones} \\
    \midrule
    Rotación de puestos              & [Sí/No/Parcial] & \\
    Ampliación de tareas             & [Sí/No/Parcial] & \\
    Enriquecimiento de puestos       & [Sí/No/Parcial] & \\
    Equipos autogestionados          & [Sí/No/Parcial] & \\
    Teletrabajo / trabajo flexible   & [Sí/No/Parcial] & \\
    Job crafting                     & [Sí/No/Parcial] & \\
    \bottomrule
  \end{tabularx}
\end{table}

\subsubsection{Rotación de puestos (\textit{job rotation})}

% TODO: ¿Existe un programa formal de rotación en CaixaBank?
% Es habitual en banca para desarrollar directivos y gestores con visión
% global. ¿Se rota entre oficinas, entre departamentos, internacionalmente
% (BPI)?

[Descripción a completar]

\textbf{Ventajas observadas en CaixaBank:}
\begin{itemize}
  \item [...]
\end{itemize}

\textbf{Inconvenientes observados:}
\begin{itemize}
  \item [...]
\end{itemize}

\subsubsection{Ampliación de tareas (\textit{job enlargement})}

% TODO: Con la digitalización bancaria, los asesores de oficina han
% asumido tareas que antes hacían otros departamentos (apertura online,
% soporte digital...). ¿Esto genera sobrecarga? ¿Hay compensación
% retributiva?

[Descripción a completar]

\subsubsection{Enriquecimiento de puestos (\textit{job enrichment})}

% TODO: ¿Se da autonomía a los asesores para tomar decisiones sobre
% productos? ¿Tienen feedback claro sobre su impacto en los resultados?
% El modelo de Hackman y Oldham es muy útil para evaluar si el diseño
% actual de los puestos genera motivación intrínseca.

[Descripción a completar]

\subsubsection{Teletrabajo y trabajo flexible}

% TODO: Tras la pandemia COVID-19, CaixaBank implementó políticas de
% teletrabajo. ¿Sigue vigente? ¿Para qué puestos? ¿Con qué porcentaje
% de jornada? ¿Ha mejorado la satisfacción y productividad?

[Descripción a completar]

\subsection{Formas de rediseño recomendadas para CaixaBank}

% TODO: Proponed al menos 2-3 formas que no se apliquen o que se apliquen
% de forma mejorable. Justificad por qué serían beneficiosas.

[Descripción a completar]

\begin{notaequipo}
  Preguntad en la entrevista:
  \begin{enumerate}
    \item ¿Existe algún programa formal de rotación? ¿Tiene criterios claros?
    \item ¿Los empleados de oficina han asumido más tareas en los últimos años?
          ¿Se ha revisado su carga de trabajo y retribución?
    \item ¿Tienen autonomía para adaptar su forma de trabajo (\textit{job crafting})?
    \item ¿Cuál es la política actual de teletrabajo? ¿Están satisfechos con ella?
    \item ¿Se mide el impacto de los rediseños en productividad y satisfacción?
  \end{enumerate}
\end{notaequipo}


% ================================================================
\section{Nuevas tendencias en análisis y diseño de puestos: el enfoque por competencias}
% ================================================================

% RECORDATORIO AL EQUIPO:
% El ADP por competencias supone un cambio de paradigma:
%
% ADP TRADICIONAL:
%   - Se centra en TAREAS y RESPONSABILIDADES del puesto.
%   - El resultado es una descripción de puesto estática.
%   - El encaje persona-puesto se hace comparando requisitos formales
%     (titulación, experiencia) con el perfil del candidato.
%
% ADP POR COMPETENCIAS:
%   - Se centra en los COMPORTAMIENTOS OBSERVABLES que llevan al éxito.
%   - Identifica competencias clave (técnicas + conductuales/transversales).
%   - Es más flexible: la misma competencia puede ser válida en varios puestos.
%   - Facilita la gestión integrada del talento bajo un marco común.
%   - Permite adaptar rápidamente los perfiles a entornos cambiantes (VUCA).
%
% OTRAS TENDENCIAS:
%   - ADP en entornos ágiles (puestos definidos por roles y sprints).
%   - Uso de People Analytics e IA para rediseñar puestos.
%   - Gestión por proyectos vs. gestión por puestos (modelos gig/fluid).
%   - ADP continuo: revisión periódica sistemática.
%   - Wellbeing y diseño de puestos orientado al bienestar.
%
% PARA ESTA SECCIÓN INDICAD:
%   - ¿Aplica CaixaBank el ADP por competencias?
%   - ¿Tiene un diccionario de competencias corporativo?
%   - ¿Cómo integra el enfoque en selección, formación y evaluación?
%   - ¿Usa People Analytics?

\subsection{Del enfoque tradicional al enfoque por competencias}

El análisis y diseño de puestos ha evolucionado desde el modelo clásico,
centrado en la descripción de tareas y responsabilidades, hacia un enfoque
basado en \textbf{competencias}, que pone el foco en los comportamientos
observables y medibles que conducen a un desempeño excelente.

\begin{table}[H]
  \centering
  \small
  \caption{Comparativa entre el ADP tradicional y el ADP por competencias}
  \begin{tabularx}{\textwidth}{p{0.23\textwidth} X X}
    \toprule
    \textbf{Dimensión}       & \textbf{ADP Tradicional} & \textbf{ADP por Competencias} \\
    \midrule
    Unidad de análisis       & Tareas y responsabilidades & Competencias y comportamientos \\
    Resultado                & Descripción de puesto estática & Perfil de competencias dinámico \\
    Orientación              & Presente                  & Presente y futuro \\
    Flexibilidad             & Baja                      & Alta \\
    Integración con RRHH     & Limitada                  & Total (marco integrado) \\
    Adaptación al cambio     & Reactiva                  & Proactiva \\
    \bottomrule
  \end{tabularx}
\end{table}

\subsection{El modelo de competencias de CaixaBank}

% TODO: CaixaBank, como gran empresa del IBEX 35, es probable que tenga
% un diccionario de competencias corporativo. Investigad:
%   - ¿Existe un modelo de competencias formalizado?
%   - ¿Cuáles son las competencias corporativas (transversales a todos)?
%   - ¿Hay competencias específicas por nivel o familia de puestos?
%   - ¿Cómo se usa el modelo en selección, evaluación y formación?

[Descripción a completar con la información de la entrevista y/o la
web corporativa de CaixaBank (sección de empleo / cultura)]

Ejemplos de competencias frecuentes en entidades bancarias
(ajustar a las reales de CaixaBank):
\begin{itemize}
  \item \textbf{Orientación al cliente}: escucha activa, empatía,
        resolución de problemas del cliente.
  \item \textbf{Orientación a resultados}: planificación, seguimiento
        de objetivos, toma de decisiones basada en datos.
  \item \textbf{Trabajo en equipo y colaboración}: comunicación, confianza,
        gestión de conflictos.
  \item \textbf{Adaptabilidad e innovación}: apertura al cambio,
        aprendizaje continuo, pensamiento creativo.
  \item \textbf{Integridad y cumplimiento}: ética profesional, respeto
        normativo, confidencialidad.
\end{itemize}

\subsubsection{Integración del modelo de competencias en los procesos de RRHH}

% TODO: ¿Cómo se usa el diccionario de competencias en:
%   - Selección (entrevista por competencias, assessment centres)?
%   - Evaluación del desempeño (evaluación 360º, por objetivos + competencias)?
%   - Formación y desarrollo (planes de desarrollo individuales)?
%   - Planes de carrera y sucesión?

[Descripción a completar]

\subsection{Otras tendencias emergentes}

\subsubsection{People Analytics aplicado al análisis de puestos}

% TODO: ¿Usa CaixaBank datos para analizar patrones de desempeño,
% rotación o satisfacción por tipo de puesto? ¿Tienen un equipo de
% People Analytics? ¿Qué sistemas usan (SAP, Workday, Cornerstone...)?

[Descripción a completar]

\subsubsection{ADP en entornos ágiles}

% TODO: Especialmente relevante en el área de tecnología y digitalización
% de CaixaBank. ¿Trabajan con metodologías ágiles (Scrum, SAFe)?
% ¿Cómo afecta eso al diseño de puestos (roles vs. puestos fijos)?

[Descripción a completar]

\subsubsection{Diseño de puestos y bienestar del empleado (\textit{wellbeing})}

% TODO: CaixaBank ha publicado informes de sostenibilidad con compromisos
% en bienestar del empleado. ¿Se refleja eso en el diseño de los puestos?
% ¿Tienen indicadores de bienestar vinculados al diseño de puestos?

[Descripción a completar]

\begin{notaequipo}
  Preguntad en la entrevista:
  \begin{enumerate}
    \item ¿Tienen un diccionario o modelo de competencias corporativo?
          ¿Podéis verlo o que os lo describan?
    \item ¿Cómo se usa ese modelo en la selección y en la evaluación
          del desempeño?
    \item ¿Tienen un equipo o función de People Analytics?
    \item ¿Trabajan con metodologías ágiles en alguna área? ¿Cómo ha
          cambiado el concepto de ``puesto'' en esas áreas?
    \item ¿Qué tendencias en ADP creen que les falta adoptar?
  \end{enumerate}
\end{notaequipo}


% ---- Capítulo 3: Tema 3 — Planificación de RRHH ----
% !TeX root = ../tfg.tex
% !TeX encoding = utf8
%
% ====================================================================
%  CAPÍTULO 3 — TEMA 3: PLANIFICACIÓN DE RECURSOS HUMANOS
%
%  Este capítulo cubre todos los apartados del Tema 3 de la asignatura:
%    3.1  Planificación de RRHH o de plantilla
%    3.2  Análisis de oferta y demanda de RRHH
%    3.3  Equilibrio entre oferta y demanda
%    3.4  Criterios de evaluación y sistemas de información (SIRH)
%    3.5  Procesos sustractivos
%    3.6  Alternativas a la reducción / incorporación de plantilla
% ====================================================================

\chapter{Tema 3: Planificación de Recursos Humanos}


% ================================================================
\section{Planificación de recursos humanos en CaixaBank}
% ================================================================

% RECORDATORIO AL EQUIPO:
% La planificación de RRHH consiste en anticipar las necesidades
% futuras de personal y definir las acciones necesarias para satisfacerlas.
%
% NIVELES DE PLANIFICACIÓN:
%   1. OPERATIVA (corto plazo, 1 año): cubrir necesidades inmediatas.
%   2. TÁCTICA (medio plazo, 1-3 años): ajuste de plantilla, formación,
%      movilidad interna.
%   3. ESTRATÉGICA (largo plazo, +3 años): alineada con la estrategia
%      de negocio; transforma la estructura de la plantilla para los
%      retos futuros (digitalización, nuevos mercados, etc.).
%
% PARA ESTA SECCIÓN INDICAD:
%   - ¿Existe un proceso formal de planificación de plantilla en CaixaBank?
%   - ¿Es operativa, táctica o estratégica (o las tres)?
%   - ¿Está integrada con el plan estratégico del grupo?
%   - ¿Está la empresa satisfecha? ¿Qué mejoraría?

La planificación de recursos humanos es el proceso mediante el cual la
organización prevé sus necesidades futuras de personal y define las
acciones necesarias ---contratación, formación, reducción,
reconversión--- para alinear la dotación de personal con los objetivos
estratégicos del negocio.

\subsection{¿Existe planificación formal de RRHH en CaixaBank?}

% TODO: CaixaBank publica planes estratégicos trienales (p.ej. Plan
% Estratégico 2022-2024). Comprobad si ese plan incluye objetivos
% específicos de gestión de personas / plantilla. La sección de
% Sostenibilidad de la Memoria Anual también incluye datos de plantilla
% y objetivos de RRHH.

[Descripción a completar con la información obtenida en la entrevista
y en los documentos públicos de CaixaBank (Memoria Anual, Plan Estratégico)]

\subsection{Carácter estratégico de la planificación de RRHH}

% Para que la planificación sea estratégica debe:
%   a) Estar integrada con el plan estratégico de negocio.
%   b) Tener un horizonte temporal de al menos 3 años.
%   c) Considerar escenarios de futuro (no solo extrapolar el pasado).
%   d) Involucrar a la alta dirección, no solo al departamento de RRHH.
%   e) Contar con indicadores de seguimiento (KPIs de plantilla).

\begin{table}[H]
  \centering
  \small
  \caption{Características de la planificación estratégica de RRHH en CaixaBank}
  \begin{tabularx}{\textwidth}{p{0.45\textwidth} p{0.12\textwidth} X}
    \toprule
    \textbf{Característica}                          & \textbf{¿Se cumple?} & \textbf{Observaciones} \\
    \midrule
    Integración con plan estratégico de negocio      & [Sí/No/Parcial] & \\
    Horizonte temporal $\geq$ 3 años                 & [Sí/No/Parcial] & \\
    Consideración de escenarios de futuro            & [Sí/No/Parcial] & \\
    Involucración de la alta dirección               & [Sí/No/Parcial] & \\
    KPIs de seguimiento definidos                    & [Sí/No/Parcial] & \\
    Planificación de la sucesión                     & [Sí/No/Parcial] & \\
    \bottomrule
  \end{tabularx}
\end{table}

\subsection{Impacto de la fusión con Bankia en la planificación de plantilla}

% La fusión CaixaBank-Bankia (2021) fue la mayor fusión bancaria de la
% historia de España. Implicó un ajuste de plantilla significativo (ERE
% 2021). Analizad cómo afectó a la planificación de RRHH: ¿fue reactiva
% o proactiva? ¿Cómo se gestionó la integración de plantillas?

[Descripción a completar]

\begin{notaequipo}
  Preguntad en la entrevista:
  \begin{enumerate}
    \item ¿Existe un plan de plantilla formal? ¿Quién lo elabora y con
          qué periodicidad?
    \item ¿Está vinculado al plan estratégico del grupo?
    \item ¿Cómo se gestionó la planificación de plantilla tras la fusión
          con Bankia?
    \item ¿Tienen planificación de sucesión para posiciones clave?
    \item ¿Qué KPIs de plantilla monitorizan (ratio empleado/oficina,
          rotación, cobertura de vacantes, etc.)?
  \end{enumerate}
\end{notaequipo}


% ================================================================
\section{Análisis de oferta y demanda de recursos humanos}
% ================================================================

% RECORDATORIO AL EQUIPO:
%
% ANÁLISIS DE DEMANDA (¿cuántos empleados necesitamos?):
%  Métodos cuantitativos:
%    - Análisis de series temporales (extrapolación de tendencias).
%    - Ratios de productividad (unidades producidas / empleado).
%    - Análisis de regresión.
%    - Modelos econométricos.
%  Métodos cualitativos:
%    - Técnica Delphi (consenso de expertos internos).
%    - Juicio directivo (managers estiman sus necesidades).
%    - Escenarios (planificación por escenarios).
%
% ANÁLISIS DE OFERTA (¿cuántos empleados tenemos / tendremos?):
%  Oferta interna (inventario de la plantilla actual):
%    - Inventario de habilidades / perfiles.
%    - Matrices de Markov / tablas de transición.
%    - Análisis de la rotación y del absentismo.
%    - Mapas de reemplazamiento / planes de sucesión.
%  Oferta externa (mercado laboral):
%    - Análisis del mercado de trabajo local / sectorial.
%    - Datos del INE, SEPE, convenios colectivos.
%    - Benchmarking con otras entidades financieras.

\subsection{Análisis de la demanda de RRHH}

La previsión de la demanda de personal permite a CaixaBank anticipar
cuántos y qué tipo de empleados necesitará en el futuro para alcanzar
sus objetivos de negocio.

\subsubsection{Métodos utilizados para la previsión de la demanda}

\begin{table}[H]
  \centering
  \small
  \caption{Métodos de previsión de la demanda de RRHH en CaixaBank}
  \begin{tabularx}{\textwidth}{p{0.30\textwidth} p{0.12\textwidth} X}
    \toprule
    \textbf{Método}                    & \textbf{¿Se usa?} & \textbf{Descripción / observaciones} \\
    \midrule
    Series temporales / tendencias     & [Sí/No/Parcial] & \\
    Ratios de productividad            & [Sí/No/Parcial] & \\
    Regresión / modelos cuantitativos  & [Sí/No/Parcial] & \\
    Técnica Delphi                     & [Sí/No/Parcial] & \\
    Juicio directivo                   & [Sí/No/Parcial] & \\
    Planificación por escenarios       & [Sí/No/Parcial] & \\
    \bottomrule
  \end{tabularx}
\end{table}

[Descripción detallada de los métodos más relevantes que usa CaixaBank]

\textbf{Métodos que serían convenientes aplicar:}

% TODO: Basándoos en el tamaño y complejidad de CaixaBank, ¿qué métodos
% más avanzados podrían añadir valor? Ejemplos:
%   - Modelos de regresión múltiple incorporando variables macroeconómicas
%     (volumen de crédito, tipos de interés, PIB).
%   - People Analytics con machine learning para predecir rotación.
%   - Planificación por escenarios ante disrupciones tecnológicas (IA).

\begin{itemize}
  \item [Método 1 — justificación]
  \item [Método 2 — justificación]
\end{itemize}

\subsection{Análisis de la oferta de RRHH}

\subsubsection{Oferta interna: inventario de la plantilla}

% TODO: ¿Tiene CaixaBank un sistema de inventario de habilidades?
% ¿Usan matrices de Markov o tablas de transición para modelar flujos
% internos? ¿Tienen mapas de reemplazamiento / planes de sucesión?

[Descripción a completar]

\subsubsection{Oferta externa: análisis del mercado de trabajo}

% TODO: ¿Cómo monitoriza CaixaBank el mercado laboral financiero?
% ¿Hacen benchmarking de salarios con competidores (Santander, BBVA)?
% ¿Colaboran con universidades para captar talento (prácticas, campus
% recruitment)?

[Descripción a completar]

\begin{notaequipo}
  Preguntad en la entrevista:
  \begin{enumerate}
    \item ¿Con qué métodos cuantifican sus necesidades futuras de personal?
    \item ¿Tienen un inventario actualizado de competencias de la plantilla?
    \item ¿Usan modelos matemáticos o estadísticos para prever la oferta
          interna (p.ej. matrices de transición / Markov)?
    \item ¿Cómo monitorizan el mercado laboral externo?
    \item ¿Cuál es su tasa de rotación actual? ¿La consideran alta o baja?
    \item ¿Han tenido dificultades para cubrir ciertos perfiles en los últimos
          años? ¿Cuáles? (Probablemente perfiles tecnológicos: Data Scientists,
          especialistas en ciberseguridad)
  \end{enumerate}
\end{notaequipo}


% ================================================================
\section{Equilibrio entre oferta y demanda de recursos humanos}
% ================================================================

% RECORDATORIO AL EQUIPO:
% Tras comparar oferta y demanda se pueden dar tres situaciones:
%
%  1. OFERTA = DEMANDA → Equilibrio. Ideal pero poco frecuente.
%     Acciones: mantenimiento, formación para cambios menores.
%
%  2. DEMANDA > OFERTA → Déficit de personal. Necesidad de incorporar.
%     Acciones: reclutamiento externo, promoción interna, horas extra,
%     trabajo temporal, outsourcing, boomerang employees, automatización.
%
%  3. OFERTA > DEMANDA → Superávit de personal.
%     Acciones: ERE, ERTE, jubilaciones anticipadas, congelación de
%     contrataciones, reducción de jornada, outplacement.
%     (Se desarrolla en la sección de procesos sustractivos)
%
% PARA ESTA SECCIÓN INDICAD:
%   - ¿Cuál es la situación de CaixaBank actualmente?
%   - ¿Ha variado en los últimos años (especialmente tras la fusión)?
%   - ¿Qué acciones concretas ha tomado para lograr el equilibrio?

\subsection{Situación actual de CaixaBank}

% TODO: Tras el ERE de 2021, la plantilla se ha ido ajustando. Pero la
% digitalización crea demanda de nuevos perfiles tecnológicos. Puede
% haber simultáneamente superávit en ciertos perfiles tradicionales
% (cajeros) y déficit en perfiles digitales.

[Descripción a completar]

\begin{table}[H]
  \centering
  \small
  \caption{Situación de equilibrio oferta-demanda en CaixaBank por tipo de perfil}
  \begin{tabularx}{\textwidth}{p{0.33\textwidth} p{0.22\textwidth} X}
    \toprule
    \textbf{Familia de puestos}    & \textbf{Situación} & \textbf{Acciones adoptadas} \\
    \midrule
    Perfiles bancarios tradicionales (oficina) & [Superávit/Déficit/Equilibrio] & \\
    Perfiles tecnológicos (datos, ciberseg., digital) & [Superávit/Déficit/Equilibrio] & \\
    Perfiles de gestión y dirección  & [Superávit/Déficit/Equilibrio] & \\
    Perfiles de cumplimiento y riesgo & [Superávit/Déficit/Equilibrio] & \\
    \bottomrule
  \end{tabularx}
\end{table}

\subsection{Acciones de ajuste ante déficit de personal}

\begin{itemize}
  \item \textbf{Reclutamiento externo:} [descripción]
  \item \textbf{Promoción interna:} [descripción]
  \item \textbf{Reskilling / upskilling:} [¿tienen programas de recualificación?]
  \item \textbf{Trabajo temporal o externalización:} [descripción]
  \item \textbf{Acuerdos con universidades / campus recruitment:} [descripción]
\end{itemize}

\subsection{Acciones de ajuste ante superávit de personal}

% Se desarrollará con mayor detalle en la sección de Procesos Sustractivos.
% Aquí haced un resumen de los mecanismos de reducción utilizados.

\begin{itemize}
  \item \textbf{ERE 2021 (fusión Bankia):} [número de afectados, condiciones]
  \item \textbf{Jubilaciones anticipadas:} [descripción]
  \item \textbf{Recolocación interna (movilidad geográfica / funcional):} [descripción]
  \item \textbf{Planes de formación para reconversión profesional:} [descripción]
\end{itemize}

\begin{notaequipo}
  Preguntad en la entrevista:
  \begin{enumerate}
    \item ¿En qué perfiles sienten más escasez de candidatos en el mercado?
    \item ¿Cómo afrontaron el exceso de plantilla tras la fusión con Bankia?
    \item ¿Han implementado programas de reskilling o upskilling para
          reconvertir perfiles tradicionales en perfiles digitales?
    \item ¿Cómo gestionan las situaciones de superávit temporal sin
          recurrir a despidos (ERTE, reducción de jornada...)?
    \item ¿Está la plantilla actual equilibrada o hay áreas con tensión?
  \end{enumerate}
\end{notaequipo}


% ================================================================
\section{Criterios de evaluación de la planificación y sistemas de información de RRHH}
% ================================================================

% RECORDATORIO AL EQUIPO:
%
% INDICADORES CLAVE DE RRHH (KPIs):
%   Indicadores de flujo de plantilla:
%     - Tasa de rotación (turnover rate): % de bajas sobre plantilla media.
%     - Tasa de absentismo: horas no trabajadas / horas contratadas.
%     - Tiempo medio de cobertura de vacantes (time-to-fill).
%     - Tasa de promoción interna.
%   Indicadores de calidad del ajuste:
%     - % de posiciones clave cubiertas por planes de sucesión.
%     - Tasa de retención en el primer año.
%   Indicadores de coste:
%     - Coste medio de contratación (cost-per-hire).
%     - ROI de formación.
%   Indicadores de previsión:
%     - Desviación entre plantilla planificada y real.
%
% SISTEMAS DE INFORMACIÓN DE RRHH (SIRH / HRIS):
%   Módulos típicos: administración de personal, gestión del tiempo,
%   selección (ATS), formación (LMS), evaluación del desempeño,
%   planificación de plantilla, People Analytics.
%   Sistemas del mercado: SAP SuccessFactors, Workday, Oracle HCM,
%   Cornerstone, Meta4 / NGA HR, A3 RRHH.

\subsection{Criterios e indicadores para evaluar la planificación de RRHH}

La eficacia de la planificación de RRHH se mide a través de indicadores
clave (KPIs) que permiten detectar desviaciones y tomar decisiones
correctoras.

\subsubsection{KPIs utilizados por CaixaBank}

\begin{table}[H]
  \centering
  \small
  \caption{Indicadores de planificación de RRHH en CaixaBank}
  \begin{tabularx}{\textwidth}{p{0.35\textwidth} p{0.15\textwidth} X}
    \toprule
    \textbf{Indicador}                    & \textbf{¿Se mide?} & \textbf{Valor / fuente} \\
    \midrule
    Tasa de rotación (turnover)           & [Sí/No] & [\% — Memoria Anual / entrevista] \\
    Tasa de absentismo                    & [Sí/No] & [\%] \\
    Tiempo medio de cobertura (TTF)       & [Sí/No] & [días] \\
    Tasa de promoción interna             & [Sí/No] & [\%] \\
    Coste medio de contratación           & [Sí/No] & [\euro{}] \\
    Desviación plantilla real vs. plan    & [Sí/No] & [\%] \\
    Cobertura de sucesión (puestos clave) & [Sí/No] & [\%] \\
    ROI de formación                      & [Sí/No] & [\%] \\
    \bottomrule
  \end{tabularx}
\end{table}

[Análisis de los indicadores más relevantes obtenidos en la entrevista]

\subsubsection{Propuestas de mejora en los criterios de evaluación}

% TODO: ¿Qué KPIs deberían medir y no miden? Por ejemplo:
%   - NPS de empleados (Employee Net Promoter Score).
%   - Índice de diversidad e inclusión.
%   - Indicadores de bienestar (wellbeing index).
%   - Métricas de People Analytics predictivo.

\begin{itemize}
  \item [Propuesta 1 — justificación]
  \item [Propuesta 2 — justificación]
\end{itemize}

\subsection{Sistemas de Información de Recursos Humanos (SIRH)}

\subsubsection{Sistema SIRH utilizado por CaixaBank}

% TODO: CaixaBank, dada su escala, es probable que use uno de los grandes
% sistemas enterprise (SAP SuccessFactors, Workday, Oracle HCM). En 2021,
% tras la fusión, habrán tenido que integrar los sistemas de ambas
% entidades. ¿Cuál es el sistema unificado resultante?

[Descripción a completar con la información de la entrevista]

\begin{table}[H]
  \centering
  \small
  \caption{Módulos del SIRH de CaixaBank}
  \begin{tabularx}{\textwidth}{p{0.42\textwidth} p{0.12\textwidth} X}
    \toprule
    \textbf{Módulo / Funcionalidad}          & \textbf{¿Activo?} & \textbf{Observaciones} \\
    \midrule
    Administración de personal / nómina       & [Sí/No] & \\
    Gestión del tiempo y asistencia           & [Sí/No] & \\
    Reclutamiento y selección (ATS)           & [Sí/No] & \\
    Gestión del desempeño                     & [Sí/No] & \\
    Formación y desarrollo (LMS)              & [Sí/No] & \\
    Planificación de plantilla                & [Sí/No] & \\
    People Analytics e informes               & [Sí/No] & \\
    Gestión de la compensación                & [Sí/No] & \\
    Portal del empleado (autoservicio)        & [Sí/No] & \\
    \bottomrule
  \end{tabularx}
\end{table}

\subsubsection{Propuestas de mejora del SIRH}

% TODO: Reflexionad sobre qué capacidades adicionales aportarían valor:
%   - Integrar People Analytics predictivo (predecir quién va a marcharse).
%   - Mejorar la experiencia del empleado en el portal de autoservicio.
%   - Integrar IA para el cribado de CV (con control del sesgo).
%   - Mejorar la interoperabilidad entre el SIRH y otros sistemas del banco.

[Descripción a completar]

\begin{notaequipo}
  Preguntad en la entrevista:
  \begin{enumerate}
    \item ¿Qué sistema SIRH utilizan? ¿Ha cambiado tras la fusión con Bankia?
    \item ¿Qué módulos tienen activos? ¿Qué módulos les gustaría activar?
    \item ¿Qué KPIs de plantilla monitorizan de forma habitual?
    \item ¿Tienen un cuadro de mando (dashboard) de RRHH? ¿Quién lo consulta?
    \item ¿Están satisfechos con la información que les proporciona el SIRH?
          ¿Qué mejorarían?
  \end{enumerate}
\end{notaequipo}


% ================================================================
\section{Procesos sustractivos: gestión del exceso de plantilla}
% ================================================================

% RECORDATORIO AL EQUIPO:
% Los procesos sustractivos son aquellos mediante los cuales la
% organización reduce su plantilla cuando la oferta de RRHH supera
% la demanda. Se clasifican en:
%
%  A) SUSPENSIÓN LABORAL TEMPORAL → ERTES
%     - El contrato se SUSPENDE temporalmente (no se extingue).
%     - El trabajador cobra prestación del SEPE + complemento empresa.
%     - Muy usado durante la pandemia COVID-19.
%     - Ventaja: mantiene la plantilla formada.
%
%  B) RUPTURA LABORAL VOLUNTARIA
%     1. Dimisión: el empleado abandona sin coste para la empresa.
%     2. Jubilación (ordinaria, anticipada o parcial).
%
%  C) RUPTURA LABORAL INVOLUNTARIA
%     1. Despido individual.
%     2. ERE — Expediente de Regulación de Empleo (despidos colectivos).
%     3. Downsizing: reducción sistemática del tamaño organizativo.
%     4. No renovación de contratos temporales.
%
%  D) RECOLOCACIÓN — OUTPLACEMENT
%     - Apoyo a los empleados que salen para que encuentren nuevo empleo.
%     - Mejora la imagen y reduce el trauma de la salida.

Cuando la oferta de RRHH supera la demanda, CaixaBank debe recurrir a
mecanismos de reducción o reajuste de plantilla. A continuación se
analizan los principales procesos sustractivos aplicados o aplicables.

\subsection{Suspensión laboral temporal: los ERTE}

% TODO: ¿Ha recurrido CaixaBank a ERTE? Principalmente durante la
% pandemia COVID-19. Si es así, ¿a cuántos empleados afectó? ¿Qué
% duración tuvo? ¿Qué complementos ofreció la empresa sobre el SEPE?

Un ERTE (Expediente de Regulación Temporal de Empleo) implica la
suspensión temporal del contrato de trabajo sin que se produzca su
extinción. El trabajador cobra una prestación por desempleo del SEPE,
que la empresa puede complementar.

\subsubsection{Aplicación en CaixaBank}

[Descripción a completar]

\textbf{Ventajas del ERTE para CaixaBank:}
\begin{itemize}
  \item [...]
\end{itemize}

\textbf{Inconvenientes / limitaciones:}
\begin{itemize}
  \item [...]
\end{itemize}

\subsection{Ruptura laboral voluntaria}

\subsubsection{Dimisión y rotación voluntaria}

% TODO: ¿Cuál es la tasa de dimisiones voluntarias en CaixaBank?
% ¿Realizan entrevistas de salida? ¿Qué factores llevan a los empleados
% a marcharse voluntariamente (mejores salarios en competencia, agotamiento...)?

[Descripción a completar]

\subsubsection{Jubilación (ordinaria y anticipada)}

% TODO: CaixaBank tiene una plantilla relativamente envejecida en
% ciertos segmentos. ¿Han implementado planes de jubilación anticipada?
% ¿En el ERE de 2021 se incluyeron medidas de jubilación anticipada?
% ¿Con qué condiciones (edad mínima, complemento sobre la pensión)?

[Descripción a completar]

\textbf{Ventajas de los planes de jubilación anticipada:}
\begin{itemize}
  \item Reducción no traumática de la plantilla sin despidos.
  \item Renovación generacional y entrada de perfiles más digitales.
  \item Menor conflicto sindical que los despidos colectivos.
\end{itemize}

\textbf{Inconvenientes:}
\begin{itemize}
  \item Pérdida de conocimiento experto y memoria institucional.
  \item Coste económico elevado para la empresa (complementos).
  \item Riesgo de que se jubilen empleados clave que no deberían marcharse.
\end{itemize}

\subsection{Ruptura laboral involuntaria: ERE y despidos}

\subsubsection{ERE 2021 — Fusión CaixaBank-Bankia}

% TODO: El ERE de 2021 fue el más grande de la historia bancaria española.
% Detalles clave a incluir:
%   - Número de afectados (aprox. 6.452 empleados según acuerdo sindical).
%   - Condiciones: indemnizaciones, bajas incentivadas, jubilaciones anticipadas.
%   - Sindicatos involucrados: UGT, CCOO, FINE.
%   - ¿Fue el ajuste suficiente o hay nuevos ajustes previstos?
%   - ¿Cómo afectó a la moral de la plantilla que se quedó?

En 2021, tras la fusión con Bankia, CaixaBank llevó a cabo el mayor
Expediente de Regulación de Empleo (ERE) de la banca española:

\begin{table}[H]
  \centering
  \small
  \caption{Datos del ERE CaixaBank 2021}
  \begin{tabularx}{\textwidth}{lX}
    \toprule
    \textbf{Parámetro}           & \textbf{Dato} \\
    \midrule
    Empleados afectados           & $\sim$6.452 (según acuerdo negociado) \\
    Cierre de oficinas            & $\sim$1.500 oficinas \\
    Tipo de salidas               & Bajas incentivadas, jubilaciones anticipadas, suspensiones \\
    Periodo de consultas          & [Fecha inicio -- Fecha acuerdo] \\
    Sindicatos firmantes          & [CCOO, UGT, FINE --- verificar] \\
    Medidas sociales de acompañamiento & [Outplacement, complementos, formación] \\
    \bottomrule
  \end{tabularx}
\end{table}

\textbf{Valoración del proceso:}

[Descripción a completar --- ¿fue percibido como justo? ¿Qué impacto
tuvo en el clima laboral? ¿Síndrome del superviviente?]

\subsubsection{Downsizing y reestructuración continua}

% TODO: Más allá del ERE de 2021, ¿hay un proceso de reducción continua
% (no renovación de contratos, digitalización que sustituye puestos)?
% La automatización de procesos bancarios reduce la necesidad de ciertos
% perfiles de back-office.

[Descripción a completar]

\subsection{Recolocación: el \textit{outplacement}}

\subsubsection{Servicios de \textit{outplacement} en el ERE de CaixaBank}

% TODO: ¿Ofreció CaixaBank servicios de outplacement a los empleados
% que salieron? ¿Lo gestionaron internamente o con consultora externa
% (Lee Hecht Harrison, Randstad RiseSmart, Adecco)?
% ¿Cuántos empleados se beneficiaron? ¿Con qué tasa de recolocación?

El outplacement es un conjunto de servicios de apoyo a los empleados
que abandonan la empresa, orientados a facilitar su reincorporación
al mercado laboral:

\begin{itemize}
  \item Orientación y coaching profesional.
  \item Actualización y elaboración del currículum.
  \item Preparación para entrevistas de trabajo.
  \item Identificación de oportunidades en el mercado.
  \item Apoyo psicológico durante la transición.
\end{itemize}

[Descripción de cómo se aplicó en CaixaBank]

\textbf{Ventajas del outplacement para CaixaBank:}
\begin{itemize}
  \item Mejora la imagen como empleador responsable (\textit{employer branding}).
  \item Reduce el conflicto y los litigios laborales.
  \item Mitiga el impacto emocional en los empleados que salen, lo que
        también reduce el síndrome del superviviente en los que se quedan.
\end{itemize}

\begin{notaequipo}
  Preguntad en la entrevista:
  \begin{enumerate}
    \item ¿Han recurrido a ERTE en algún momento (COVID, fusión)?
    \item ¿Qué condiciones se ofrecieron en el ERE de 2021?
          ¿Hubo jubilaciones anticipadas? ¿A partir de qué edad?
    \item ¿Se ofreció outplacement a los empleados afectados por el ERE?
          ¿Quién lo gestionó? ¿Cuáles fueron los resultados?
    \item ¿Cómo afectó el ERE al clima laboral de los que se quedaron?
          ¿Tomaron medidas para gestionar ese impacto?
    \item ¿Realizan entrevistas de salida cuando un empleado dimite?
          ¿Qué hacen con esa información?
  \end{enumerate}
\end{notaequipo}


% ================================================================
\section{Alternativas a la reducción de plantilla e incorporación a la empresa}
% ================================================================

% RECORDATORIO AL EQUIPO:
%
% ALTERNATIVAS A LA REDUCCIÓN (cuando hay superávit):
%   1. Reducción de la jornada laboral (trabajo a tiempo parcial).
%   2. Reparto del trabajo (job sharing / worksharing).
%   3. Reducción salarial temporal (sacrificio compartido).
%   4. Congelación de contrataciones y de salarios.
%   5. Formación y recualificación (reskilling).
%   6. Movilidad funcional: asignar empleados a otras áreas con demanda.
%   7. Movilidad geográfica: trasladar empleados a zonas con necesidades.
%   8. Excedencias voluntarias y licencias sin sueldo.
%   9. Reducción de horas extra.
%  10. Externalización de funciones no estratégicas (outsourcing).
%
% ALTERNATIVAS A LA INCORPORACIÓN (cuando hay déficit):
%   1. Promoción interna (movilidad vertical).
%   2. Formación acelerada para cubrir el gap de habilidades.
%   3. Horas extra (si el déficit es temporal y limitado).
%   4. Trabajo temporal o contratos de proyecto.
%   5. Externalización (BPO, outsourcing, contratas).
%   6. Automatización (RPA, IA) para reducir la necesidad de personal.
%   7. Programas de boomerang employees (recontratación de exempleados).

\subsection{Alternativas a la reducción de plantilla}

Antes de recurrir a medidas de ruptura laboral, CaixaBank puede ---y
debe--- explorar alternativas menos traumáticas para ajustar la
plantilla preservando el empleo y el clima organizacional.

\begin{table}[H]
  \centering
  \small
  \caption{Alternativas a la reducción de plantilla en CaixaBank}
  \begin{tabularx}{\textwidth}{p{0.37\textwidth} p{0.12\textwidth} X}
    \toprule
    \textbf{Alternativa}                        & \textbf{¿Se aplica?} & \textbf{Observaciones} \\
    \midrule
    Reducción de jornada / trabajo parcial       & [Sí/No] & \\
    Congelación de contrataciones                & [Sí/No] & \\
    Reskilling / reconversión de perfiles        & [Sí/No] & \\
    Movilidad funcional interna                  & [Sí/No] & \\
    Movilidad geográfica                         & [Sí/No] & \\
    Excedencias voluntarias incentivadas         & [Sí/No] & \\
    Outplacement previo al despido               & [Sí/No] & \\
    Reducción de horas extra / subcontratación   & [Sí/No] & \\
    \bottomrule
  \end{tabularx}
\end{table}

\subsubsection{Reskilling y recualificación profesional}

% TODO: Este es probablemente el aspecto más relevante para CaixaBank.
% La digitalización bancaria hace que muchos perfiles de red de oficinas
% queden obsoletos. ¿Tiene CaixaBank un programa de reskilling para
% reconvertir a estos empleados en perfiles digitales?
% Mencionad si existe un programa concreto (p.ej. ``Future Skills'' o similar).

[Descripción a completar]

\textbf{Ventajas:}
\begin{itemize}
  \item Preserva el empleo y el conocimiento del negocio.
  \item Mejor imagen como empresa responsable.
  \item Menor coste que despedir y contratar externamente.
\end{itemize}

\textbf{Inconvenientes:}
\begin{itemize}
  \item No todos los empleados tienen el potencial o la motivación para
        reconvertirse en perfiles muy técnicos.
  \item Requiere inversión elevada en formación.
  \item El proceso es lento; no resuelve necesidades inmediatas de
        perfiles digitales especializados.
\end{itemize}

\subsubsection{Movilidad interna (funcional y geográfica)}

% TODO: ¿Tiene CaixaBank un sistema formal de movilidad interna?
% ¿Los empleados pueden postularse a vacantes internas? ¿Hay un portal
% de oportunidades internas en el SIRH? ¿Con qué frecuencia se producen
% traslados voluntarios o forzosos?

[Descripción a completar]

\subsection{Alternativas a la incorporación de nueva plantilla}

Cuando existe déficit de personal, CaixaBank no siempre recurre de
inmediato a la contratación externa. Las alternativas disponibles
incluyen:

\subsubsection{Promoción interna}

% TODO: ¿Cuál es la política de CaixaBank respecto a la promoción
% interna? ¿Se publican primero las vacantes internamente antes de
% salir al mercado? ¿Tienen planes de carrera definidos?

[Descripción a completar]

\subsubsection{Automatización y digitalización}

% TODO: La RPA y la IA están transformando el sector bancario. CaixaBank
% ha invertido fuertemente en tecnología. ¿Hay procesos que antes
% requerían empleados y ahora están automatizados? ¿Esto ha reducido
% la necesidad de incorporaciones en ciertas áreas?

[Descripción a completar]

\subsubsection{Externalización (\textit{outsourcing})}

% TODO: ¿Qué funciones tiene CaixaBank externalizadas?
% ¿IT, call centers, servicios de back-office, seguridad, limpieza?
% ¿Han seguido la tendencia de insourcing en algún área?

[Descripción a completar]

\subsubsection{Trabajo temporal y contratos de proyecto}

[Descripción a completar]

\begin{notaequipo}
  Preguntad en la entrevista:
  \begin{enumerate}
    \item ¿Tienen algún programa de reskilling o upskilling para reconvertir
          perfiles tradicionales? ¿En qué consiste? ¿Cuántos empleados se
          han beneficiado?
    \item ¿Existe un portal de vacantes internas para la movilidad interna?
    \item ¿Cómo deciden entre cubrir una vacante internamente o acudir
          al mercado externo?
    \item ¿Qué procesos han automatizado recientemente con RPA o IA?
          ¿Ha supuesto una reducción de la plantilla en esas áreas?
    \item ¿Qué funciones tienen externalizadas? ¿Han internalizado alguna
          función en los últimos años?
  \end{enumerate}
\end{notaequipo}


% -------------------------------------------------------------------
%  APPENDIX
%  (apéndices, numerados con letras)
% -------------------------------------------------------------------
\appendix
\pdfbookmark[-1]{Apéndices}{appendix}

% Apéndice A: Preguntas y respuestas de la entrevista
% !TeX root = ../tfg.tex
% !TeX encoding = utf8
%
% ====================================================================
%  APÉNDICE A — PREGUNTAS REALIZADAS A LA EMPRESA Y RESPUESTAS
%
%  RECORDATORIO AL EQUIPO:
%  - Este apéndice es OBLIGATORIO según las instrucciones de la asignatura.
%  - Debéis incluir las preguntas que realizasteis en la entrevista/visita
%    a la empresa y las respuestas que os dieron.
%  - Organizad las preguntas por temática (Tema 2 y Tema 3) o por el
%    orden en que se realizaron.
%  - Podéis incluir el nombre o cargo del entrevistado (si da permiso)
%    y la fecha y medio de la entrevista (presencial, videollamada...).
%  - Si hay información que el entrevistado pidió que no se publicara,
%    indicadlo (p.ej. "Dato confidencial — no se puede reproducir").
%  - Si grabasteis la entrevista (con permiso), podéis transcribir los
%    fragmentos más relevantes.
% ====================================================================

\chapter{Entrevista a CaixaBank: Temas 2 y 3 (análisis de puestos y planificación)}

% ============================================================
\section{Datos de la entrevista}
% ============================================================

\begin{table}[H]
  \centering
  \begin{tabularx}{\textwidth}{lX}
    \toprule
    \textbf{Parámetro}       & \textbf{Detalle} \\
    \midrule
    Empresa                  & CaixaBank \\
    Nombre del entrevistado  & [Anónimo a petición] \\
    Cargo                    & [Director de Recursos Humanos regional] \\
    Departamento             & [Recursos Humanos] \\
    Fecha                    & [13/03/2026] \\
    Duración                 & [15 min] \\
    Formato                  & [Telefónica] \\
    \bottomrule
  \end{tabularx}
\end{table}

% ============================================================
\section{Preguntas sobre el Tema 2 — Análisis y diseño de puestos}
% ============================================================

%Jose Angel
\subsubsection*{En clase hemos estudiado que el Análisis del Puesto de Trabajo es el proceso a través del cual la empresa recopila y analiza la información sobre
 el puesto para identificar sus tareas y responsabilidades. En una entidad tan grande como CaixaBank, cuando necesitáis definir un puesto
(por ejemplo, de nueva creación o uno que ha cambiado mucho), ¿cómo es vuestro proceso paso a paso para analizarlo? ¿Sentís que este proceso os funciona
 bien o es demasiado burocrático?}

\textbf{Respuesta:}\newline
``Los puestos los define un departamento Secretaria Tecnica de Negocios. Deloitte desarrollo el modelo. De abajo arriba, desde las necesidades del cliente. No se hace desde recursos humanos la definición, esta externalizado. Se basan mucho en las soft skills (dependiendo del lugar donde está la oficina).''
%David

\subsubsection*{Sabemos que existen algunos trabajos que pueden resultar monotonos o terminar quemando al empleado. ¿Qué medidas tomáis para rediseñar puestos monótonos o, por el contrario, muy especializados? ¿Habéis aplicado técnicas como la rotación de puestos, la ampliación de tareas o dar más autonomía a los empleados?}

\textbf{Respuesta:}\newline
``Desde la empresa se trabaja desde dos ámbitos: Una forma reactiva, de forma que se publican las vacantes frecuentemente y cada trabajador puede postular al que este interesado; y otra forma proactiva, en la que se hacen entrevistas individuales para ver si un empleado concreto encaja con el puesto para el cual estamos buscando empleado.''

\subsubsection*{En CaixaBank, a la hora de definir lo que necesitáis de un trabajador, ¿os centráis más en una lista estricta de tareas o vais más hacia un modelo de competencias blandas y adaptabilidad?}

\textbf{Respuesta:}\newline
``Los puestos se definen principalmente basandose en las tareas que van er, de forma que nos centramos principalmente en las "Hard-skills. Existe un departamento técnico específico que define las necesidades de cada puesto.''

% ============================================================
\section{Preguntas sobre el Tema 3 — Planificación de RRHH}
% ============================================================

% JESÚS 
\subsubsection*{Pasando a un plano más estratégico, la planificación de recursos humanos trata de anticiparse y prever el movimiento de personas. ¿Cómo planificáis las necesidades de personal a futuro? ¿Dirías que es una mera planificación administrativa de cuánta gente falta o está totalmente alineada con la estrategia de negocio del banco, es decir, pensáis en las necesidades a corto plazo (jubilaciones, etc.) o lo planificáis más a largo plazo?}

\textbf{Respuesta:}\newline
``En el tema de dimensionamineto se tienen en cuenta muchísimos factores y es un aspecto clave al determinar toda la estructura de costes. Los factores más importante son el mercado actual y potencial de cada plaza (vacante o puesto), la su estacionalidad o la tipología de clientes (por ejemplo, suelen tener distintas necesidades en pueblos que en ciudades) entre otros que son fundamentales para algo muy importante para nosotros como es el dimensionamiento.''

\subsubsection*{¿Y hay alguna herramienta o método concreto que useis para estas previones de demanda u oferta de trabajadores?}

\textbf{Respuesta:}\newline
``Se está trabajando con IA y, como la inmensa mayoría de las empresas, en todo el tema de planificación se usa SAP, que es un programa de gestión (SIRH) desde hace muchísimos años que ha evolucionado y que se usa en empresas desde el ámbito de la construcción hasta recursos humanos (casi todas las grandes empresas lo usan para eso). Respecto a la IA tenemos licencias de Copilot.''

% JULIAN ----------------------------------------------------------------------------------------------------
% - ¿Hay Sistemas de Información de RRHH? ¿En qué consisten? ¿Cómo pueden mejorar? 
% - Procesos Sustractivos: DEMANDA MENOR A OFERTA 
%   o ¿Existe suspensión laboral (excedencias/ERTES)? 
%   o ¿Existe ruptura laboral voluntaria (dimisión, jubilación, jubilación anticipada)? 

\subsubsection*{Para manejar a tantos miles de empleados, imaginamos que la tecnología es vital. Los sistemas de información de recursos humanos (SIRH) recopilan, registran y almacenan los datos de los empleados. ¿Qué tipo de plataforma utilizáis en CaixaBank para esto? ¿Estáis contentos con ella o creéis que le falta evolucionar en algún aspecto?”}

\textbf{Respuesta:}\newline
``Ahora mismo se esta trabajando con IA, y bueno, nosotros como la inmensa mayoría de empresas, utilizamos una herramienta llamada SAP que se utiliza en todos los procesos de planificación. SAP es un programa de gestión desde hace muchísimos años. Todas las grandes empresas trabajamos en el entorno SAP. Y ahora como te digo, también herramientas de IA, en concreto tenemos como herramienta de cabecera y licenciada Copilot. ''

\subsubsection*{Hablando del movimiento de personas, a veces los trabajadores pausan su actividad, interrumpiendo temporalmente la prestación laboral. ¿Cómo gestionáis temas como las excedencias voluntarias o por cuidado de familiares? ¿Es fácil reincorporar el talento después?}
\textbf{Respuesta:}\newline
`` Siempre que el empleado se preste (en la mayoría de los casos lo hace), realizamos entrevistas de salida y tratamos de organizar cual ha sido el punto de valor que ha provocado la ruptura. ''

\subsubsection*{Y cuando alguien decide irse definitivamente, ya sea por dimisión o por jubilación (normal o anticipada), ¿cómo gestionáis esas salidas? ¿Hacéis entrevistas de salida para entender los motivos y mejorar la retención del talento?}
\textbf{Respuesta:}\newline
`` Para la retención de talento, aunque a veces se toma por medidas económicas (donde tenemos una libertad limitada), intentamos no mantener exclusivamente por motivos económicos. Lo que trabajamos mucho es en el desarrollo profesional dentro de la empresa. Trabajamos mucho el tema de la formación, sobretodo en puestos técnicos donde hay una mayor tasa de abandono. Desarrollamos un plan de carrera en base a expectativas, y el tema de oportunidades, dado que nuestra empresa tiene una capacidad que muy pocas empresas en este país tienen.''

% ------------------PARTE DE ISMAEL-----------------

\subsubsection*{En ocasiones, por estrategia o economía, es la empresa quien debe finalizar la relación laboral. En un sector de tanta transformación como la banca, que ha sufrido reestructuraciones y reducciones de plantilla, ¿cómo se intenta gestionar un proceso tan delicado como un despido colectivo (ERE) o salidas individuales para que el impacto humano sea el menor posible?}

\textbf{Respuesta:}\newline
``En colectivos queda acotado por el sindicato. En individuales (disciplinarios), ateniéndose a la ley, si se produce no hay forma de mantener al empleado.''

\subsubsection*{Antes de llegar a reducir plantilla, ¿se suelen explorar alternativas dentro del banco, como reasignaciones a otros departamentos, puestos compartidos o cambios en las políticas de empleo? ¿Creéis que hay margen para incorporar nuevas medidas?}

\textbf{Respuesta:} \newline
``Se intenta en la medida de lo posible, pero no suele ocurrir.''

\subsubsection*{Por último, hemos estudiado que los programas de recolocación u outplacement ayudan mucho a los empleados que salen a encontrar un nuevo rumbo profesional. ¿CaixaBank ofrece este tipo de servicios o apoyo emocional y formativo para facilitar su integración en el mercado laboral? ¿Qué tal es la experiencia con esto?}

\textbf{Respuesta:}\newline
``No se trabaja, pero se ha trabajado en época por ejemplo de la fusión. Proporcionándoles la entrada en empresas del mismo grupo.''



% -------------------------------------------------------------------
%  GLOSARIO (opcional — comentad si no lo usáis)
% -------------------------------------------------------------------
% % !TeX root = tfg.tex
% !TeX encoding = utf8
%
% ====================================================================
%  GLOSARIO DE TÉRMINOS (OPCIONAL)
%  Activadlo descomentando % !TeX root = tfg.tex
% !TeX encoding = utf8
%
% ====================================================================
%  GLOSARIO DE TÉRMINOS (OPCIONAL)
%  Activadlo descomentando % !TeX root = tfg.tex
% !TeX encoding = utf8
%
% ====================================================================
%  GLOSARIO DE TÉRMINOS (OPCIONAL)
%  Activadlo descomentando \input{glosario} en tfg.tex.
% ====================================================================

\chapter*{Glosario}
\addcontentsline{toc}{chapter}{Glosario}

\begin{description}
  \item[APT] Análisis de Puestos de Trabajo.
  \item[ADP] Análisis y Diseño de Puestos.
  \item[ERE] Expediente de Regulación de Empleo.
  \item[ERTE] Expediente de Regulación Temporal de Empleo.
  \item[KPI] \textit{Key Performance Indicator}, indicador clave de rendimiento.
  \item[RRHH] Recursos Humanos.
  \item[SIRH] Sistema de Información de Recursos Humanos.
  \item[TTF] \textit{Time-to-Fill}, tiempo medio de cobertura de una vacante.
\end{description}

\endinput
 en tfg.tex.
% ====================================================================

\chapter*{Glosario}
\addcontentsline{toc}{chapter}{Glosario}

\begin{description}
  \item[APT] Análisis de Puestos de Trabajo.
  \item[ADP] Análisis y Diseño de Puestos.
  \item[ERE] Expediente de Regulación de Empleo.
  \item[ERTE] Expediente de Regulación Temporal de Empleo.
  \item[KPI] \textit{Key Performance Indicator}, indicador clave de rendimiento.
  \item[RRHH] Recursos Humanos.
  \item[SIRH] Sistema de Información de Recursos Humanos.
  \item[TTF] \textit{Time-to-Fill}, tiempo medio de cobertura de una vacante.
\end{description}

\endinput
 en tfg.tex.
% ====================================================================

\chapter*{Glosario}
\addcontentsline{toc}{chapter}{Glosario}

\begin{description}
  \item[APT] Análisis de Puestos de Trabajo.
  \item[ADP] Análisis y Diseño de Puestos.
  \item[ERE] Expediente de Regulación de Empleo.
  \item[ERTE] Expediente de Regulación Temporal de Empleo.
  \item[KPI] \textit{Key Performance Indicator}, indicador clave de rendimiento.
  \item[RRHH] Recursos Humanos.
  \item[SIRH] Sistema de Información de Recursos Humanos.
  \item[TTF] \textit{Time-to-Fill}, tiempo medio de cobertura de una vacante.
\end{description}

\endinput


% -------------------------------------------------------------------
%  BACKMATTER
%  (bibliografía)
% -------------------------------------------------------------------
\backmatter
\pdfbookmark[-1]{Referencias}{BM-Referencias}

\bibliographystyle{alpha-es}
% Otros estilos: abbrv, plain, alpha-es, plain-es
\begin{small}
  \bibliography{library.bib}
\end{small}

\end{document}
